% Complete standalone source. Compile with pdfLaTeX (two passes).
\documentclass[11pt,letterpaper]{article}
\usepackage[T1]{fontenc}
\usepackage{lmodern}
\usepackage{amsmath,amssymb,amsthm,mathtools}
\usepackage[margin=1in]{geometry}
\usepackage{microtype}
\usepackage{booktabs,tabularx,array}
\usepackage{enumitem}
\usepackage[authoryear,round]{natbib}
\usepackage{xcolor}
\usepackage{hyperref}
\usepackage{bookmark}
\usepackage{fancyhdr}
\definecolor{linkblue}{RGB}{22,62,100}
\hypersetup{colorlinks=true,linkcolor=linkblue,citecolor=linkblue,urlcolor=linkblue,
 pdfauthor={Colombo Team},
 pdftitle={Unrestricted Randomized Complexity of Finite-Sum Nonconvex-Strongly-Concave Minimax Optimization},
 pdfsubject={Oracle complexity under individual component smoothness},
 pdfkeywords={minimax optimization, finite sums, randomized lower bounds, individual smoothness}}
\urlstyle{same}
\newtheorem{theorem}{Theorem}[section]
\newtheorem{lemma}[theorem]{Lemma}
\newtheorem{proposition}[theorem]{Proposition}
\newtheorem{corollary}[theorem]{Corollary}
\theoremstyle{definition}
\newtheorem{definition}[theorem]{Definition}
\theoremstyle{remark}
\newtheorem{remark}[theorem]{Remark}
\newcommand{\R}{\mathbb R}
\newcommand{\E}{\mathbb E}
\newcommand{\PP}{\mathbb P}
\newcommand{\eps}{\varepsilon}
\newcommand{\cF}{\mathcal F}
\newcommand{\cR}{\mathcal R}
\newcommand{\T}{\mathfrak T}
\newcommand{\norm}[1]{\left\lVert#1\right\rVert}
\newcommand{\ip}[2]{\left\langle#1,#2\right\rangle}
\DeclareMathOperator*{\argmax}{arg\,max}
\DeclareMathOperator{\range}{range}
\DeclareMathOperator{\spanop}{span}
\newcolumntype{Y}{>{\raggedright\arraybackslash}X}
\setlist{nosep,leftmargin=*}
\setlength{\emergencystretch}{1.5em}
\setlength{\headheight}{14pt}
\setlength{\parskip}{0.2em}
\pagestyle{fancy}
\fancyhf{}
\fancyhead[L]{\small Unrestricted randomized finite-sum minimax complexity}
\fancyfoot[C]{\thepage}
\renewcommand{\headrulewidth}{0.3pt}
\renewcommand{\footrulewidth}{0pt}
\title{\bfseries Unrestricted Randomized Complexity of Finite-Sum\\
 Nonconvex--Strongly-Concave Minimax Optimization}
\author{Colombo Team\\[3pt]
 \small\href{mailto:colombo.2026@outlook.com}{\texttt{colombo.2026@outlook.com}}}
\date{}

\begin{document}
\maketitle
\thispagestyle{plain}
\begin{abstract}
We study the randomized incremental first-order complexity of smooth finite-sum
nonconvex--strongly-concave minimax optimization on unconstrained Euclidean
domains. With $n$ components, smoothness $L$, average dual strong concavity
$\mu$, initial primal gap $\Delta$, and $\kappa=L/\mu$, we establish the lower bound
\[
 \Omega\!\left(n+\frac{L\Delta}{\varepsilon^2}
 \min\{\sqrt n\,\kappa,\ n^{3/4}\sqrt\kappa\}\right)
\]
under individual component smoothness. The bound holds against unrestricted
randomized adaptive algorithms, even when a query returns the exact component
value and full gradient. Together with existing upper bounds under
mean-squared averaged smoothness, it characterizes both classes up to
logarithmic factors, under exact dual initialization and a universal
small-accuracy condition. The construction separates information recovery
from curvature calibration. A fixed distribution of orthogonally branched
linear maps remains hard under arbitrary forward and transpose probes.
Quadratic saddle relays turn this information cost into primal interactions:
a long relay covers large condition numbers, while a one-edge relay uses a
public scalar correction to tune the entire eliminated energy. A globally
smooth primal activation mechanism then forces sequential source recovery.
The result covers expected query cost with expected squared-gradient accuracy
and fixed query budgets with constant success probability. The worst case is
taken over finite dimensions, and strong concavity is required only of the
average objective; the componentwise-concave individually smooth subclass
remains outside the characterization.
\end{abstract}

\noindent\textbf{Keywords:} minimax optimization; finite-sum optimization;
oracle complexity; randomized lower bounds; individual smoothness;
averaged smoothness.

\begin{samepage}
\noindent\textbf{AI Usage.}
Nearly the entire research pipeline for this paper was carried out by
Colombo, powered by GPT-5.6 Sol.
\par
\end{samepage}

\section{Introduction}
\label{sec:introduction}

Consider the finite-sum minimax problem
\begin{equation}
 \min_{x\in\R^{d_x}}\max_{y\in\R^{d_y}} f(x,y),
 \qquad f(x,y)=\frac1n\sum_{i=1}^n f_i(x,y),
 \qquad \Phi(x)=\max_y f(x,y).
 \label{eq:problem}
\end{equation}
The average objective is strongly concave in $y$ and may be nonconvex in $x$.
Strong concavity makes the value function $\Phi$ differentiable, so a natural
goal is to find an $x$ with small $\|\nabla\Phi(x)\|$. The finite-sum structure
allows algorithms to query individual components and reuse information
across iterations. The resulting complexity depends jointly on the number
of components and the conditioning of the inner maximization.

Let $L$ denote smoothness, $\mu$ average dual strong concavity, and $\Delta$
an upper bound on the initial primal gap. Write
\begin{equation}
 A=\frac{L\Delta}{\varepsilon^2},\qquad
 \kappa=\frac L\mu,\qquad
 \cR(n,\kappa)=\min\{\sqrt n\,\kappa,\ n^{3/4}\sqrt\kappa\}.
 \label{eq:rate-notation}
\end{equation}
The polynomial rate $n+A\cR(n,\kappa)$ is now understood under
mean-squared averaged smoothness. \citet{zhou2026} establish it for randomized
zero-respecting first-order methods, and \citet{pengprox2026} establish an
unrestricted randomized lower bound even with exact component proximal
access. A distinct question remains: does the same difficulty persist when
\emph{every} component has smoothness constant at most the same numerical $L$?

This regularity distinction matters. Mean-squared averaged smoothness permits
individual component constants as large as $\sqrt n\,L$. A construction
that concentrates an interaction in one component may use precisely this
extra curvature. Renormalizing such a construction to make each component
$L$-smooth can change both the condition number and the initial gap, and
therefore change the resulting lower bound. The unrestricted algorithm model
adds a separate difficulty: an algorithm may query an arbitrary dense
direction, use an arbitrarily large norm, or output a point it has never
queried. An exact oracle can reveal information through an arbitrarily small
nonzero response.

We prove that individual smoothness still permits the full rate
$\Omega(n+A\cR(n,\kappa))$ against unrestricted randomized algorithms.
Existing averaged-smooth algorithms give the matching upper bound, up to
logarithms. The theorem applies for every $n\ge1$ and $\kappa\ge1$ under
a universal condition $A\ge A_*$; it imposes no additional relation between
accuracy, $n$, and $\kappa$.

\paragraph{Main contributions.}
The contribution is an individually smooth lower-bound construction with
three interacting parts.
\begin{enumerate}[label=(\roman*)]
 \item \textbf{A fixed information law for arbitrary probes.}
 Hidden isometries are distributed across orthogonal component branches.
 A chronological accounting argument bounds source recovery even when each
 round supplies forward and transpose actions on every layer and the oracle
 may reveal intermediate sources. The instance distribution is fixed before
 the algorithm's random seed is drawn.
 \item \textbf{Two compatible curvature calibrations.}
 A long quadratic saddle relay gives the $n^{3/4}\sqrt\kappa$ regime.
 A one-edge relay covers the $\sqrt n\,\kappa$ regime by cancelling an
 order-one endpoint self-energy with a public scalar quadratic. Both
 constructions retain uniform individual smoothness.
 \item \textbf{A global primal stationarity certificate.}
 All nonlinear activation terms depend only on primal variables. The
 components are globally smooth on the entire primal--dual space, and a
 small value-function gradient forces recovery of the terminal source.
 The reduction includes random stopping and unqueried outputs.
\end{enumerate}

The hard instances require strong concavity of the \emph{average} dual
objective. Individual dual quadratics may be indefinite. This property is
part of the model and is essential to the construction; the result does not
settle the smaller individually smooth class with componentwise dual
concavity.

\section{Model and main theorem}
\label{sec:model}

All spaces are finite-dimensional Euclidean spaces. Norms of vectors are
Euclidean, and norms of linear or multilinear maps are the induced operator
norms. The dimension is known to the algorithm and is unrestricted in the
worst case.

\subsection{Function classes and initialization}

Each $f_i:\R^{d_x+d_y}\to\R$ is continuously differentiable. We distinguish
two smoothness conditions, using the same numerical parameter $L$:
\begin{align}
 \|\nabla f_i(z)-\nabla f_i(z')\|
 &\le L\|z-z'\| \quad (i\in[n]),
 &&\text{individual smoothness (IS)},\label{eq:IS}\\
 \frac1n\sum_{i=1}^n\|\nabla f_i(z)-\nabla f_i(z')\|^2
 &\le L^2\|z-z'\|^2,
 &&\text{averaged smoothness (AS)}.\label{eq:AS}
\end{align}
Here $z=(x,y)$ and $[n]=\{1,\ldots,n\}$. IS implies AS; AS implies that
the average $f$ is $L$-smooth by Jensen's inequality. Smoothness of $f$
alone is weaker than AS.

For every $x$, assume that $f(x,\cdot)$ is $\mu$-strongly concave, with
$0<\mu\le L$. Its maximizer
$y^*(x)=\argmax_y f(x,y)$ is unique, and
$\nabla\Phi(x)=\nabla_x f(x,y^*(x))$. We impose
\begin{equation}
 \Phi(0)-\inf_x\Phi(x)\le\Delta,
 \qquad y^*(0)=0.
 \label{eq:initialization}
\end{equation}
The second condition is \emph{exact dual initialization}. There is no bound
on query norms, component gradient magnitudes, or $\|y^*(x)\|$ away from the
initial point. We denote the two classes by
$\cF^0_{\mathrm{IS}}(n,L,\mu,\Delta)$ and
$\cF^0_{\mathrm{AS}}(n,L,\mu,\Delta)$. The parameters $L$ and $\mu$ are
bounds and need not be attained.

\subsection{Oracle access and performance criteria}

An enriched incremental first-order oracle returns
\begin{equation}
 \mathcal O_f(i,x,y)
 =\bigl(f_i(x,y),\nabla_x f_i(x,y),\nabla_y f_i(x,y)\bigr).
 \label{eq:oracle}
\end{equation}
The usual gradient-only oracle omits the value. A component query costs one
unit and a full pass costs $n$ units. An algorithm chooses component indices,
points, stopping times, and its output using arbitrary measurable functions
of its transcript and an independent private random seed. Computation and
memory are free. Its first query is unrestricted, and its output need not be
a previously queried point. The instance is fixed independently of the seed.

For $s\in\{\mathrm{IS},\mathrm{AS}\}$, let $\T_s^{\mathrm E}$ be the infimum
of the worst-case expected number of queries among algorithms satisfying
\begin{equation}
 \E\|\nabla\Phi(\widehat x)\|^2\le\varepsilon^2
 \quad\text{on every instance in }\cF_s^0.
 \label{eq:expected-criterion}
\end{equation}
Let $\T_s^{\mathrm P}$ be the smallest deterministic query budget guaranteeing
\begin{equation}
 \PP\bigl(\|\nabla\Phi(\widehat x)\|\le\varepsilon\bigr)\ge\frac23
 \quad\text{on every instance in }\cF_s^0.
 \label{eq:probability-criterion}
\end{equation}
These criteria are treated separately: truncating a random-time algorithm
controls its success probability, but need not preserve an expected
squared-gradient guarantee for a fallback output.

\begin{theorem}[Unrestricted randomized complexity]
\label{thm:main}
There are universal constants $A_*,c>0$ such that, for every integer $n\ge1$,
every $L,\Delta,\varepsilon>0$, and every $\mu\in(0,L]$ with
$A=L\Delta/\varepsilon^2\ge A_*$,
\begin{equation}
 c\bigl[n+A\cR(n,\kappa)\bigr]
 \le \T^{\mathrm E}_{\mathrm{IS}}
 \le \T^{\mathrm E}_{\mathrm{AS}}
 \le \widetilde O\bigl(n+A\cR(n,\kappa)\bigr).
 \label{eq:main-complexity}
\end{equation}
The same inequalities hold with $\T^{\mathrm E}$ replaced by
$\T^{\mathrm P}$. The lower bounds hold for the enriched oracle
\eqref{eq:oracle}, and the upper bounds use only gradients. The lower bound
has no logarithmic loss. All constants are independent of $n,L,\mu,\Delta,
\varepsilon$ and the dimension.
\end{theorem}

The two branches meet at $\kappa=\sqrt n$. Equivalently, for either
smoothness class and either performance criterion, the complexity is
\begin{equation}
 \begin{cases}
 \widetilde\Theta\!\left(n+\dfrac{L\Delta}{\varepsilon^2}
                 \sqrt n\,\kappa\right),
       &1\le\kappa\le\sqrt n,\\[5pt]
 \widetilde\Theta\!\left(n+\dfrac{L\Delta}{\varepsilon^2}
                 n^{3/4}\sqrt\kappa\right),
       &\kappa\ge\sqrt n.
 \end{cases}
 \label{eq:two-regimes}
\end{equation}
The notation $\widetilde O$ and $\widetilde\Theta$ suppresses the
logarithmic factors in the cited upper algorithms. The theorem concerns
worst-case complexity over finite dimensions, rather than a dimension bound
uniform in the other parameters. The universal accuracy threshold is not
optimized.

\paragraph{Structure of the hard instances.}
Both condition-number-dependent constructions have the form
\begin{equation}
 f_i(x,y)=h(x)+\langle Cx,y\rangle-\frac12\langle y,H_i y\rangle,
 \qquad \frac1n\sum_{i=1}^n H_i\succeq\mu I,
 \label{eq:bilinear-structure}
\end{equation}
with common $h$ and common $C$. Component-dependent nonlinear
primal--dual coupling is therefore unnecessary. The component dual matrices
$H_i$ may be indefinite. The separate minimization and component-search
families used for bounded $\kappa$ and for the additive $n$ term need not
share the common-$h$ representation.

\section{Comparison with prior work}
\label{sec:related}

\paragraph{Rates, regularity, and admissible algorithms.}
\citet{zhang2021} established finite-sum NC--SC lower bounds for linear-span
methods and a Catalyst upper bound. \citet{han2024} developed a
partitioned-chain construction for a restricted proximal incremental
first-order protocol. Table~\ref{tab:lower-comparison} aligns the polynomial
rates, parameter ranges, regularity assumptions, and algorithm classes.
For the earlier results, it displays a common sufficient range rather than
their largest individual ranges.

\begin{table}[t]
\centering
\small
\setlength{\tabcolsep}{3.5pt}
\renewcommand{\arraystretch}{1.18}
\begin{tabularx}{\textwidth}{@{}p{0.20\textwidth}p{0.19\textwidth}p{0.11\textwidth}p{0.08\textwidth}Y@{}}
\toprule
Reference & Lower bound & Range & Class & Algorithms and replies\\
\midrule
\citet{zhang2021}
 & $\Omega(n+A\sqrt{n\kappa})$ & $\kappa\ge2n$ & AS
 & Linear-span; component gradients\\
\citet{han2024}
 & $\Omega(n+A\sqrt{n\kappa})$ & $\kappa\ge2n$ & AS
 & Restricted PIFO protocol; values, gradients, and proximal replies\\
\citet{zhou2026}
 & $\Omega(n+A\cR)$ & All $\kappa\ge1$ & AS
 & Randomized zero-respecting; component gradients\\
\citet{pengprox2026}
 & $\Omega(n+A\cR)$ & All $\kappa\ge1$ & AS
 & Unrestricted randomized; values, gradients, and exact component proximal replies\\
Theorem~\ref{thm:main}
 & $\Omega(n+A\cR)$ & All $\kappa\ge1$ & IS
 & Unrestricted randomized; values and full component gradients\\
\bottomrule
\end{tabularx}
\caption{Lower-bound comparison, with $A=L\Delta/\varepsilon^2$ and
$\cR=\min\{\sqrt n\,\kappa,n^{3/4}\sqrt\kappa\}$.
For a uniform comparison, take $n\ge2$ and
$\varepsilon^2\le cL\Delta$ for a sufficiently small universal $c$.
The first two rows use sufficient common ranges; the Han et al. row takes
the primal weak-convexity parameter to be $L$.
AS and IS use the same numerical $L$. Zhou's result and the two
unrestricted lower bounds admit exact dual initialization; the earlier rows
retain their original initialization conventions. Each oracle and algorithm
restriction is part of the result.}
\label{tab:lower-comparison}
\end{table}

The unrestricted AS result of \citet{pengprox2026} is particularly relevant:
its hard family has polynomial dimension and individually strongly concave
dual components, but its individual smoothness constants may scale as
$\sqrt n\,L$. Our theorem establishes the same rate under IS, using a
larger-dimensional construction and average dual strong concavity. The two
results strengthen different aspects of the model. In particular, our
value-and-gradient lower bound does not imply a component-proximal lower
bound. Their upper bound allows inexact dual initialization but assumes
$\|\nabla f(0,0)\|^2\le2L\Delta$; the matching upper bounds used here
assume $y^*(0)=0$.

\paragraph{A componentwise-concave bilinear subclass.}
The IS bilinear comparison in \citet{zhou2026} assumes
$f_i(x,y)=g_i(x)+\langle A_i x,y\rangle-h_i(y)$ with each $h_i$
$\mu$-strongly convex. Its lower bound is
$\Omega(n+A\min\{\kappa,\sqrt{n\kappa}\})$ for zero-respecting methods.
The separable algorithms of \citet{jin2022} exploit additional structure
of this kind. Our construction generally lies outside that subclass;
Table~\ref{tab:lower-comparison} therefore concerns the general NC--SC class.

\paragraph{Information hiding and saddle constructions.}
Random rotations are a standard tool for unrestricted stationary-point
lower bounds \citep{carmon2020}. For finite sums,
\citet{zhougu2019} study structured linear-span lower bounds, while
\citet{pengdense2026} give an unrestricted individually smooth minimization
benchmark. We use the latter only for bounded condition numbers.
Our information problem instead uses orthogonal component branches and
explicitly accounts for source disclosures and transpose probes under one
fixed product-Haar law. Dual-path constructions for NC--SC problems are
developed in \citet{li2021}. Here the endpoint normalization and the
one-edge cancellation are tailored to individual smoothness and to a
single global primal stationarity certificate.

\paragraph{Algorithmic inputs.}
The upper bound is supplied by existing methods, as summarized in
Table~\ref{tab:upper-regimes}. The present work contributes the IS lower
bound and its unrestricted information reduction.

\begin{table}[t]
\centering
\small
\renewcommand{\arraystretch}{1.22}
\begin{tabularx}{\textwidth}{@{}p{0.19\textwidth}p{0.25\textwidth}Y@{}}
\toprule
Condition number & Matching complexity & Gradient-only upper algorithm\\
\midrule
$1\le\kappa\le\sqrt n$
 & $\widetilde\Theta(n+A\sqrt n\,\kappa)$
 & RPDA \citep[Theorem 3]{zhou2026}\\
$\kappa\ge\sqrt n$
 & $\widetilde\Theta(n+An^{3/4}\sqrt\kappa)$
 & Catalyst with variance reduction \citep[Theorem 4.1, Corollary 4.1, and Section 4.2]{zhang2021}\\
\bottomrule
\end{tabularx}
\caption{The two regimes in Theorem~\ref{thm:main}. Both upper algorithms
apply under AS, average dual strong concavity, and exact dual initialization.
The lower bounds hold under IS. For $n=1$, an upper algorithm for two
identical components changes the cost only by a universal factor.}
\label{tab:upper-regimes}
\end{table}

\section{How the lower bound is assembled}
\label{sec:overview}

The proof has three interfaces: a source-recovery problem, a quadratic
realization of source transport, and an activation mechanism that makes
small primal gradients require successful transport. This section explains
their roles and the resulting rate. The appendices give the full
construction and proofs.

\subsection{Orthogonal branches and chronological information cost}

At depth $t$, let $\mathcal H_t=(\R^n)^{\otimes t}\otimes\R^D$.
The canonical branch map is $J_{t,i}a=e_i\otimes a$. Independent rotations
$O_t$ define
\begin{equation}
 A_{t,i}=O_{t+1}J_{t,i}O_t^{\mathsf T},\qquad
 U_t=\frac1{\sqrt n}\sum_i A_{t,i},\qquad
 u_t=O_t(s^{\otimes t}\otimes e),
 \label{eq:overview-source}
\end{equation}
where $s=n^{-1/2}(1,\ldots,1)$ and $O_0$ fixes the public unit vector $e$.
The branches have orthogonal ranges, $U_t$ is an isometry, and
$U_tu_t=u_{t+1}$. Thus all $n$ branches contribute equally to the next
source direction.

We analyze an oracle stronger than the component oracle: one round can
request a forward and a transpose action on every edge, using different
component indices across edges. All inputs precede all replies. The oracle
may also disclose entire intermediate sources at designated sites. A source
at site $j$ is eligible only once the represented fraction $p_j$ reaches
a fixed threshold $\rho$; disclosure then fills the accounting sets through
depth $j$. If $p_t(K)$ is the represented fraction after $K$ rounds and
$r_K$ is the number of disclosures, the key bound is
\begin{equation}
 p_t(K)\le\min\left\{1,\rho^{-r_K}\frac{\binom Kt}{n^t}\right\}.
 \label{eq:overview-mass}
\end{equation}
A disclosure increases represented mass by at most a factor $\rho^{-1}$.
This charge preserves information cached before a stage becomes active.

Arbitrary dense inputs are handled by an exact conditional-Haar exposure.
Each new input requires at most one fresh Gaussian direction in the analysis. The last-factor
slices of all these directions span a subspace $V\subset\R^D$ with
$\E\|P_Ve\|^2=\dim(V)/D$. Every represented unit vector $z$ at depth $t$
satisfies
\begin{equation}
 |\langle u_t,z\rangle|^2
 \le p_t+(1-p_t)\|P_Ve\|^2.
 \label{eq:overview-correlation}
\end{equation}
Choosing a sufficiently large finite $D$ makes the second term small.
The resulting serial lower bound is $\Omega(nH)$ for a path of $H$ edges,
including arbitrary query norms, private randomness, and an unqueried
output. Appendix~\ref{app:source} proves this statement without choosing a
resisting instance after the algorithm's seed.

\subsection{Quadratic relays with a common endpoint geometry}

The information maps enter a component only through dual bilinear terms.
After maximizing the average over the dual variables, both relay families
produce
\begin{equation}
 \theta Q(a,b)
 =\theta\bigl(\|a\|^2-\langle b,Pa\rangle+\langle b,Db\rangle\bigr),
 \quad \frac15I\preceq D\preceq I,\quad DP=P,
 \label{eq:overview-endpoint}
\end{equation}
where $P$ is the product of the average isometries along the relay.
The same endpoint geometry is essential: it permits one stationarity
argument to cover both regimes, including directions orthogonal to the
hidden source.

\paragraph{The long relay.}
Let $\beta=n^{-1/2}$ and take $N\ge2$ vertices. The negative average dual
Hessian is $\beta\mathcal M_N$, where
\[
 \langle w,\mathcal M_Nw\rangle
 =N^{-2}\sum_{t=0}^{N-1}\|w_t\|^2
   +\sum_{t=0}^{N-2}\|w_{t+1}-U_tw_t\|^2.
\]
The smallest eigenvalue of $\mathcal M_N$ is $N^{-2}$, so the average dual
curvature is $\beta/N^2$. Public endpoint coefficients normalize
the eliminated energy to \eqref{eq:overview-endpoint} with $\theta=1$.
Each raw component is globally $8$-smooth, even though the average
curvature is $\beta/N^2$. The proof in
Theorem~\ref{thm:long-relay} decomposes the entire terminal space into
source and innovation subspaces.

\paragraph{The one-edge relay.}
For a prescribed $\mu_0\in(0,1]$, define
\begin{equation}
 t=\mu_0+\beta,\qquad \theta=\frac\beta t,\qquad
 r^2=t-\frac{\beta^2}{t},\qquad d=\theta-\frac12.
 \label{eq:overview-short-parameters}
\end{equation}
With primal connectors $a,b$ and dual variables $w_0,w_1$, set
\begin{equation}
\begin{split}
 \mathcal J_i(a,b;w_0,w_1)
 ={}&d(\|a\|^2+\|b\|^2)
       +r(\langle a,w_0\rangle-\langle b,w_1\rangle)\\
 &-\frac t2(\|w_0\|^2+\|w_1\|^2)
       +\langle w_1,A_iw_0\rangle.
\end{split}
\label{eq:overview-short-relay}
\end{equation}
The average is $\mu_0$-strongly concave, every component is globally
$8$-smooth, and its value is \eqref{eq:overview-endpoint} with
\[
 P=U,\qquad D=I-\frac\theta2(I-UU^{\mathsf T}).
\]
The scalar $d$ cancels the order-one endpoint energy created by dual
elimination, leaving the whole interaction at amplitude $\theta$.
The hidden projector $UU^{\mathsf T}$ appears only in the eliminated
value, never in the component oracle. Negative $d$ is admissible because
primal convexity is not required. Lemma~\ref{lem:short-relay} gives the exact
Schur complement calculation.

\subsection{From stationarity to sequential recovery}

Place $M$ relays in series, with $m$ edges per relay. Stage $j$ has source
$v_j=u_{jm}$, state $S_j$, and connectors $a_j,b_j$. A smooth bounded radial
map $Z$ defines the score $t_j=\langle v_j,Z(S_j)\rangle$. Gates with exactly
flat inactive regions activate a stage only after its predecessor is
recovered. The exit uses the public vector inner product
$\langle Z(S_j),E_\nu(b_j)\rangle$, where $E_\nu$ is another radial cap.
All nonlinear terms are purely primal.

The explicit normalized components satisfy, with universal constants
$L_*,C_\Delta,\tau_*>0$,
\begin{equation}
\begin{gathered}
 \operatorname{Lip}(\nabla f_i^0)\le L_*,\qquad
 y^*(0)=0,\qquad \Phi_0(0)=0,\\
 \Phi_0(0)-\inf\Phi_0\le C_\Delta\theta M,\qquad
 \|\nabla\Phi_0(x)\|\le\theta\tau_*
       \ \Longrightarrow\ t_M\ge1.
\end{gathered}
\label{eq:overview-interface}
\end{equation}
The first line follows from local derivative bounds and bounded incidence
between adjacent stages. For the last implication, the proof first
localizes every unfinished state, then excludes intermediate scores and a
low score followed by a high score. There is consequently at most one
unfinished frontier. Connector coercivity makes the radial caps inactive
there, and the projected connector equations force a state derivative
bounded away from zero. Appendices~\ref{sec:construction}
and~\ref{app:stationarity} give the constants and estimates.

Before a source is disclosed, $q,h,V$ and their first derivatives vanish
at its score on the common hiding event; complementary gates reduce to
public constants. A component value
and full gradient can therefore be simulated exactly using one forward and
one transpose action per edge. A small-gradient output reveals a vector
correlated with $v_M$. The source theorem then gives
\begin{equation}
 \Omega(nMm)
 \label{eq:overview-information-cost}
\end{equation}
component queries. This product follows from the chronological theorem
for the full serialized construction.

\subsection{Scaling and the two regimes}

Scale variables by $s>0$ and function values by $\alpha=Ls^2/L_*$.
Taking
\[
 s=\frac{4L_*\varepsilon}{L\theta\tau_*}
\]
makes the normalized certificate threshold correspond to $4\varepsilon$.
The physical gap is at most a universal constant times
$(\varepsilon^2/L)M/\theta$. Thus $M=\Theta(A\theta)$ stages fit when
$A\theta$ exceeds a universal constant, and
\eqref{eq:overview-information-cost} becomes $\Omega(An\theta m)$.
The calibrations are
\begin{equation}
\begin{array}{c|c|c|c}
 \text{relay} & \theta & m & An\theta m\\ \hline
 \text{one edge} & \asymp \kappa/\sqrt n & 1 & \asymp A\sqrt n\,\kappa\\
 \text{long path} & 1 & \asymp\sqrt\kappa/n^{1/4}
                  & \asymp An^{3/4}\sqrt\kappa.
\end{array}
\label{eq:overview-calibration}
\end{equation}
These comparisons apply in the respective regimes up to universal
constant-width overlaps. The exact short-relay choice is
$\theta=\kappa/(L_*\sqrt n+\kappa)$.
When $A\theta$ is below the integer-calibration threshold,
$A\cR(n,\kappa)=O(n)$, so an explicit
hidden-component search supplies the required bound. Bounded condition
numbers use the minimization benchmark. Appendices~\ref{app:calibration}
and~\ref{app:protocols} cover all integer and boundary cases without an
$n$- or $\kappa$-dependent accuracy restriction.

\section{Upper bounds and performance protocols}
\label{sec:upper}

For $n\ge2$ and $\kappa\le\sqrt n$, RPDA has a fixed query budget
$\widetilde O(n+A\sqrt n\,\kappa)$ and guarantees
\eqref{eq:expected-criterion} \citep[Theorem 3]{zhou2026}. For
$\kappa\ge\sqrt n$, Catalyst with a variance-reduced inner solver has
expected cost $\widetilde O(A(n+n^{3/4}\sqrt\kappa))$ with the same
accuracy criterion \citep[Theorem 4.1, Corollary 4.1, and Section 4.2]{zhang2021}.
Since $n^{3/4}\sqrt\kappa\ge n$ in this range, these costs give the upper
bound in Theorem~\ref{thm:main}.
For $n=1$, duplicate the single component as in
Section~\ref{sec:upper-interfaces}.

Both algorithms require AS and strong concavity of the average. Exact dual
initialization removes their initial dual-error term. The relevant variance
estimate, for a fresh uniform index $I$ and history-measurable points
$z,z'$, is
\[
 \E_I\|\nabla f_I(z)-\nabla f_I(z')
           -[\nabla f(z)-\nabla f(z')]\|^2
 \le L^2\|z-z'\|^2.
\]
It follows directly from AS and does not replace $L$ by $\sqrt n\,L$.

\paragraph{Expected query cost.}
Below a fixed budget $B\asymp nMm$, the scaled hard distribution allows
success at gradient threshold $4\varepsilon$ with probability at most
$1/16$. An algorithm satisfying \eqref{eq:expected-criterion} has success
probability at least $15/16$ at that threshold. Truncating an algorithm of
worst-case expected cost $Q$ after $B$ calls loses at most $Q/B$ in success
probability, and hence $Q\ge7B/8$. Appendix~\ref{sec:expected-lower}
handles sub-unit budgets explicitly.

\paragraph{Fixed budgets and constant confidence.}
Run an upper algorithm at accuracy $\varepsilon/4$, and stop it after
$\lceil12Q_0\rceil$ calls, where $Q_0$ is its public expected-cost bound.
The probabilities of runtime overflow and of an inaccurate untruncated
output are at most $1/12$ and $1/16$, respectively. The success probability
therefore exceeds $2/3$, with the same asymptotic cost. This proves the
fixed-budget statement without asserting an expected-square guarantee for
the truncated fallback.

\section{Discussion}
\label{sec:discussion}

\paragraph{Where individual nonconcavity enters.}
For the one-edge relay, the negative component dual Hessian is
\[
 H_i=\begin{pmatrix}tI&-A_i^{\mathsf T}\\-A_i&tI\end{pmatrix}.
\]
Its selected-branch eigenvalues are $t-1$ and $t+1$. When $t<1$, a component
is not dual-concave, whereas the average has minimum eigenvalue
$t-1/\sqrt n=\mu_0>0$. Individual smoothness bounds the magnitude of the
Hessian and permits this sign pattern. Extending the same rate to the
individually smooth componentwise-concave class would require a different
construction or a sharper algorithm for that class.

\paragraph{Dimension and initialization.}
The source construction has finite dimension that can grow exponentially
with the serialized path length. The theorem is therefore a query-complexity
characterization over dimensions, and does not establish a fixed-dimensional
or arithmetic-runtime lower bound. Exact dual initialization is built into
every hard instance. Removing it preserves the lower bound, but an upper
bound may then require an additional initialization term.

\paragraph{Scope of the characterization.}
The result concerns exact incremental first-order information and direct
stationarity of $\Phi$. It makes no claim for component-proximal,
higher-order, finite-precision, or fresh-stochastic oracles. All
lower-construction formulas, simulations, and parameter choices are given
below; only the identified minimization benchmark and upper algorithms are
imported. The main conclusion is that, for the general average-strongly-concave
class, requiring individual component smoothness does not improve the
worst-case polynomial rate available under AS.


\clearpage
\begin{thebibliography}{99}
\setlength{\bibsep}{5pt}
\bibitem[Carmon et~al.(2020)Carmon, Duchi, Hinder, and Sidford]{carmon2020}
Yair Carmon, John~C. Duchi, Oliver Hinder, and Aaron Sidford.
\newblock Lower bounds for finding stationary points I.
\newblock \emph{Mathematical Programming}, 184:71--120, 2020.
\newblock \url{https://doi.org/10.1007/s10107-019-01406-y}.

\bibitem[Han et~al.(2024)Han, Xie, and Zhang]{han2024}
Yuze Han, Guangzeng Xie, and Zhihua Zhang.
\newblock Lower complexity bounds of finite-sum optimization problems: The results and construction.
\newblock \emph{Journal of Machine Learning Research}, 25(2):1--86, 2024.
\newblock \url{https://jmlr.org/papers/v25/21-0264.html}.

\bibitem[Jin et~al.(2022)Jin, Sidford, and Tian]{jin2022}
Yujia Jin, Aaron Sidford, and Kevin Tian.
\newblock Sharper rates for separable minimax and finite sum optimization via primal-dual extragradient methods.
\newblock In \emph{Proceedings of the 35th Conference on Learning Theory}, volume~178 of \emph{Proceedings of Machine Learning Research}, pages 4362--4415, 2022.
\newblock \url{https://proceedings.mlr.press/v178/jin22b.html}.

\bibitem[Li et~al.(2021)Li, Tian, Zhang, and Jadbabaie]{li2021}
Haochuan Li, Yi Tian, Jingzhao Zhang, and Ali Jadbabaie.
\newblock Complexity lower bounds for nonconvex-strongly-concave min-max optimization.
\newblock In \emph{Advances in Neural Information Processing Systems}, volume~34, 2021.
\newblock \url{https://arxiv.org/abs/2104.08708}.

\bibitem[Peng et~al.(2026a)Peng, Tang, and Jia]{pengdense2026}
Yuxing Peng, Zhiqing Tang, and Weijia Jia.
\newblock Dense weak hiding: Closing complexity gaps in nonconvex and {PL} finite-sum optimization under individual smoothness.
\newblock \emph{arXiv preprint arXiv:2609.00045v3}, 2026a.
\newblock \url{https://arxiv.org/abs/2609.00045v3}.

\bibitem[Peng et~al.(2026b)Peng, Tang, and Jia]{pengprox2026}
Yuxing Peng, Zhiqing Tang, and Weijia Jia.
\newblock On the complexity of finite-sum nonconvex-strongly-concave minimax optimization with exact proximal oracles.
\newblock \emph{arXiv preprint arXiv:2611.04806v2}, 2026b.
\newblock \url{https://arxiv.org/abs/2611.04806v2}.

\bibitem[Zhang et~al.(2021)Zhang, Yang, Guzm\'an, Kiyavash, and He]{zhang2021}
Siqi Zhang, Junchi Yang, Crist\'obal Guzm\'an, Negar Kiyavash, and Niao He.
\newblock The complexity of nonconvex-strongly-concave minimax optimization.
\newblock In \emph{Proceedings of the 37th Conference on Uncertainty in Artificial Intelligence}, volume~161 of \emph{Proceedings of Machine Learning Research}, pages 482--492, 2021.
\newblock \url{https://proceedings.mlr.press/v161/zhang21c.html}.

\bibitem[Zhou(2026)]{zhou2026}
Qihao Zhou.
\newblock Near-optimal complexity of finite-sum nonconvex-strongly-concave minimax optimization.
\newblock \emph{arXiv preprint arXiv:2610.04944v1}, 2026.
\newblock \url{https://arxiv.org/abs/2610.04944}.

\bibitem[Zhou and Gu(2019)]{zhougu2019}
Dongruo Zhou and Quanquan Gu.
\newblock Lower bounds for smooth nonconvex finite-sum optimization.
\newblock In \emph{Proceedings of the 36th International Conference on Machine Learning}, volume~97 of \emph{Proceedings of Machine Learning Research}, pages 7574--7583, 2019.
\newblock \url{https://proceedings.mlr.press/v97/zhou19b.html}.
\end{thebibliography}

\clearpage
\appendix
\section*{Appendices: complete construction and proofs}
\addcontentsline{toc}{section}{Appendices: complete construction and proofs}
The appendices follow the logical order of the lower bound.
Appendix~\ref{app:source} establishes the fixed-law source problem.
Appendix~\ref{sec:relays} realizes the two quadratic relays, and
Appendix~\ref{sec:construction} defines the global primal activation terms.
Appendix~\ref{app:stationarity} proves the gradient certificate and the
exact oracle simulation. Appendix~\ref{app:calibration} chooses the
parameters, and Appendix~\ref{app:protocols} completes the stopping-time,
baseline, and upper-bound arguments.

\section{A fixed-distribution source-recovery bound}
\label{app:source}

This appendix separates the information argument from the optimization
construction.  We define a distribution of hidden isometries and prove that
transporting a known source through $H$ layers requires order $nH$ oracle
rounds.  The oracle is deliberately stronger than an incremental gradient
oracle: every round can query every layer, transpose actions are available,
and selected intermediate sources can be disclosed.  Accounting for those
disclosures is essential when the result is used to serialize optimization
stages.

\subsection{Orthogonal branches and the enlarged oracle}
\label{app:source:model}

Fix integers $n,H\geq 1$ and a multiplicity $D\geq 1$.  For
$0\leq t\leq H$, let
\[
  \mathcal H_t=(\mathbb R^n)^{\otimes t}\otimes\mathbb R^D,
  \qquad d_t=n^tD.
\]
Let $e\in\mathbb R^D$ be a fixed unit vector, put
$s=n^{-1/2}(1,\ldots,1)$, and set $\xi_t=s^{\otimes t}$, with
$\xi_0=1$.  For $i\in[n]$, define the canonical branch maps
\[
  J_{t,i}a=e_i\otimes a,
  \qquad
  J_{t,s}=\frac1{\sqrt n}\sum_{i=1}^n J_{t,i}.
\]
They satisfy $J_{t,i}^{\mathsf T}J_{t,j}=\mathbf 1_{\{i=j\}}I$,
and $J_{t,s}$ is an isometry.

Sample $O_t$ independently from Haar measure on the orthogonal group
$O(d_t)$ for $1\leq t\leq H$.  Independently, sample $O_0$ from Haar
measure on the subgroup of $O(D)$ fixing $e$.  The entire instance is
sampled before interaction.  Define
\begin{equation}
  \begin{split}
    A_{t,i}&=O_{t+1}J_{t,i}O_t^{\mathsf T},
    \qquad U_t=O_{t+1}J_{t,s}O_t^{\mathsf T},\\
    u_t&=O_t(\xi_t\otimes e).
  \end{split}
  \label{eq:source-law}
\end{equation}
Thus $u_0=e$, $U_tu_t=u_{t+1}$, and
\[
  A_{t,i}^{\mathsf T}A_{t,j}=\mathbf 1_{\{i=j\}}I,
  \qquad \|A_{t,i}\|_{\mathrm{op}}=1,
  \qquad \frac1n\sum_{i=1}^n A_{t,i}=\frac1{\sqrt n}U_t.
\]

\paragraph{Matrix queries.}
In round $k$, the algorithm chooses an index $i_{t,k}\in[n]$ and
vectors $a_t\in\mathcal H_t$, $b_t\in\mathcal H_{t+1}$ for every
edge $0\leq t<H$.  All inputs must be chosen before any reply from
that round is returned.  The replies are
\begin{equation}
  \bigl(A_{t,i_{t,k}}a_t,\ A_{t,i_{t,k}}^{\mathsf T}b_t\bigr)_{0\leq t<H}.
  \label{eq:source-round}
\end{equation}
Indices may differ between edges.  Query norms, adaptive computation,
memory, and the final output are otherwise unrestricted.

\paragraph{Intermediate source disclosures.}
Fix $a\in(0,1)$ and write $\rho=a^2/2$.  Let
$\mathcal S\subseteq\{1,\ldots,H\}$ be a public set with at most
$M\leq H$ sites.  Each source at a site in $\mathcal S$ can be
disclosed once.  We track sets of words
$\Lambda_t\subseteq[n]^t$, initialized by
$\Lambda_0=\{\varnothing\}$ and $\Lambda_t=\varnothing$ for $t>0$.
A matrix round updates all sets simultaneously according to
\begin{equation}
  \widetilde\Lambda_0=\{\varnothing\},\qquad
  \widetilde\Lambda_t=
  \Lambda_t\cup
  \{(i_{t-1,k},w):w\in\Lambda_{t-1}\},\quad t\geq1.
  \label{eq:source-word-update}
\end{equation}
Here the right-hand sides use the sets from before the round.  Define
the represented mass at depth $t$ by
$p_t=|\Lambda_t|/n^t$.  After the matrix update, an undisclosed source
$u_j$ with $j\in\mathcal S$ is eligible for disclosure if $p_j\geq\rho$.
A disclosure returns the entire physical vector $u_j$ and resets
the accounting sets to
\begin{equation}
  \Lambda_t\leftarrow[n]^t,\qquad 0\leq t\leq j.
  \label{eq:source-reset}
\end{equation}
This reset grants additional accounting support; it does not reveal
bases or rotations at earlier layers.  Eligible disclosures can be
performed in any order, including immediately and exhaustively after
each round.  The sets, indices, and disclosure flags may be public.
Eligibility depends only on the available history and the accounting
sets, so it never conditions on an unexposed part of a rotation.

\subsection{Chronological accounting}
\label{app:source:counting}

\begin{lemma}[Deletion closure and the cost of a reset]
\label{lem:source-reset}
At every time, $(i,w)\in\Lambda_{t+1}$ implies
$w\in\Lambda_t$.  Consequently,
$p_0\geq p_1\geq\cdots\geq p_H$.
A legal disclosure increases each coordinate of the mass vector by
at most a factor $\rho^{-1}$.
\end{lemma}

\begin{proof}
Deletion closure holds initially.  Every word inserted by
\eqref{eq:source-word-update} has a suffix in the old lower set,
which remains in the new lower set.  Previously present words keep
their suffixes.  A reset fills all layers through its disclosure
site and therefore preserves closure, including at the interface
with the next layer.

There are at most $n$ possible words with any given suffix, whence
$|\Lambda_{t+1}|\leq n|\Lambda_t|$ and $p_{t+1}\leq p_t$.
If site $j$ is disclosed, then $p_t\geq p_j\geq\rho$ for all
$t\leq j$.  Replacing these entries by one multiplies each by at
most $\rho^{-1}$; the other entries remain unchanged.
\end{proof}

\begin{theorem}[Reachable mass under adaptive disclosures]
\label{thm:source-counting}
After $K$ matrix rounds and $r_K$ source disclosures,
\begin{equation}
  p_t(K)\leq
  \min\left\{1,\ \rho^{-r_K}\frac{\binom Kt}{n^t}\right\},
  \qquad 0\leq t\leq H.
  \label{eq:source-mass}
\end{equation}
The inequality holds pathwise, for every adaptive choice of indices
and every allowed disclosure schedule.  We use $\binom Kt=0$ for
$t>K$.
\end{theorem}

\begin{proof}
Before the disclosures following a round, the simultaneous update
satisfies
\[
  \widetilde p_t\leq p_t+\frac1n p_{t-1},\qquad t\geq1.
\]
Induct on rounds and insert each disclosure as a separate operation.
The update preserves the claimed bound by Pascal's identity:
\[
  \frac{\binom{k-1}{t}}{n^t}
  +\frac1n\frac{\binom{k-1}{t-1}}{n^{t-1}}
  =\frac{\binom kt}{n^t}.
\]
For $t=0$, the right-hand side of the untruncated bound is at least
one.  Each disclosure contributes at most one factor $\rho^{-1}$
by Lemma~\ref{lem:source-reset}.  Finally, $p_t\leq1$ gives
\eqref{eq:source-mass}.  The argument uses no probabilistic
independence.
\end{proof}

The disclosure factor cannot simply be omitted: revealing $u_j$
can insert all $n^j$ source coordinates at layer $j$.  The factor
$\rho^{-r_K}$ charges for that information while retaining all
probes cached before disclosure.  Since there are at most $M\leq H$
disclosures, the propagation deficit in the binomial term will
dominate their total contribution.

\subsection{Exact exposure of the fixed rotations}
\label{app:source:coupling}

The counting result controls source directions generated by branch
actions.  Arbitrary dense queries require a second argument.  We
expose the already sampled rotations only as needed and place every
new direction in a common Gaussian subspace.  Formal coordinates
in this proof are an analysis device; they are not extra information
given to the algorithm.

\begin{lemma}[Lazy Haar exposure with source disclosures]
\label{lem:source-lazy}
Fix a deterministic algorithm, a budget of $K$ matrix rounds, an
allowed disclosure schedule, and a terminal output.  The experiment
admits a coupling with mutually independent standard Gaussian vectors
\[
  g_{t,j}\sim N(0,I_{d_t}),\qquad
  0\leq t\leq H,\quad 1\leq j\leq2K+1,
\]
with the following properties.  The formal coordinate $O_t^{\mathsf T}x$
of each new physical input lies in the span of previously exposed
formal vectors at depth $t$ and one fresh raw Gaussian $g_{t,j}$.
Matrix replies have exactly their canonical $J_{t,i}$ or
$J_{t,i}^{\mathsf T}$ coordinates.  A disclosed source has formal
coordinate $\xi_t\otimes e$ and uses no Gaussian slot.  One
additional slot covers the final output.  The instance marginal is
the fixed product-Haar law in \eqref{eq:source-law}.
\end{lemma}

\begin{proof}
At vertex $t$, maintain an exposed formal subspace $K_t$, an
exposed physical subspace $L_t$, and an isometry
$V_t:K_t\to L_t$ specifying the exposed restriction of $O_t$.
Initially only $e\mapsto e$ at the root is prescribed.
Conditional on this augmented exposure history, the unexposed
completions are independent uniform isometries between the
corresponding orthogonal complements.

Consider a physical input $x$.  Its known part has formal coordinate
$V_t^*P_{L_t}x$.  If $x_\perp=(I-P_{L_t})x\neq0$, take the next
independent raw Gaussian $g$ and put
\[
  h=\frac{(I-P_{K_t})g}{\|(I-P_{K_t})g\|},
  \qquad
  O_t^{\mathsf T}x=V_t^*P_{L_t}x+\|x_\perp\|h.
\]
Declare $O_th=x_\perp/\|x_\perp\|$.  This is precisely the
conditional Haar law of the inverse image of a fixed complementary
physical direction.  The new formal vector lies in
$\operatorname{span}(K_t,g)$.  If $x_\perp=0$, no new direction is
needed.

For a matrix reply, its formal coordinate is already determined by
the appropriate canonical branch action.  Map its component
orthogonal to $K_t$ to a uniform physical direction orthogonal to
$L_t$.  This is the conditional Haar forward-image law.  Process a
source disclosure in the same way, using the specified formal vector
$\xi_t\otimes e$.  The disclosure decision is measurable with
respect to the exposure history and accounting sets, so it imposes
no further condition on an unexposed completion.  Each update fixes
only an additional restriction of one rotation; the remaining
completions preserve the stated conditional independence.

Process all inputs in a round before its replies.  There are at
most two physical inputs per vertex per round: one for the forward
action on the outgoing edge and one for the transpose action on
the incoming edge.  Process the terminal output by one more
inverse-image step.  Generate all unused Gaussian slots independently.
Finally complete the remaining partial rotations according to their
conditional Haar laws.  The resulting joint experiment has exactly
the distribution obtained by sampling all rotations before
interaction.
\end{proof}

Identify each $g_{t,j}$ with $n^t$ vectors in $\mathbb R^D$ by
taking slices in its last tensor factor.  Define
\begin{equation}
  S_H=\sum_{t=0}^H n^t,\qquad
  m_G=(2K+1)S_H,\qquad
  V=\operatorname{span}\{\text{all last-factor slices of all }g_{t,j}\}.
  \label{eq:source-gaussian-space}
\end{equation}
Including unused slots is useful: these are exactly $m_G$
independent standard Gaussian vectors, without adaptive selection.
For $D>m_G$, rotational invariance and almost-sure rank $m_G$ imply
\begin{equation}
  \mathbb E\|P_Ve\|^2=\frac{m_G}{D}.
  \label{eq:source-gaussian-projection}
\end{equation}
Indeed, the expected projector is $(m_G/D)I$.

\begin{lemma}[Represented space and source correlation]
\label{lem:source-correlation}
Every exposed formal vector at depth $t$ lies in
\begin{equation}
  \mathcal R_t=
  \bigl(E_t\otimes\operatorname{span}\{e\}\bigr)
  +\bigl((\mathbb R^n)^{\otimes t}\otimes V\bigr),
  \qquad E_t=\operatorname{span}\{e_w:w\in\Lambda_t\}.
  \label{eq:source-represented-space}
\end{equation}
This includes inputs before a matrix reply, all replies, disclosed
sources, and the final output, using the current accounting sets and
the common Gaussian space $V$.  If $O_t^{\mathsf T}z\in\mathcal R_t$
and $\|z\|\leq1$, then
\begin{equation}
  |\langle u_t,z\rangle|^2
  \leq p_t+(1-p_t)\|P_Ve\|^2
  \leq p_t+\|P_Ve\|^2.
  \label{eq:source-correlation}
\end{equation}
\end{lemma}

\begin{proof}
Write $\mathcal N_t=(\mathbb R^n)^{\otimes t}\otimes V$.
Every raw Gaussian belongs to its corresponding $\mathcal N_t$,
and canonical forward and transpose maps carry these spaces into
one another.  Lemma~\ref{lem:source-lazy} therefore places a new
arbitrary input in the represented space.  A forward reply inserts
one branch index, exactly as recorded in
\eqref{eq:source-word-update}.  A transpose reply deletes a branch
index from an old source word; deletion closure places the result
in the existing lower set.  A disclosed source is covered by the
full support granted at its site in \eqref{eq:source-reset}.
Induction over the exposure steps proves
\eqref{eq:source-represented-space}.  The extra Gaussian slots cover
unqueried outputs.

Put $v=\|P_Ve\|^2$.  Split the last tensor factor into $V$ and
$V^\perp$.  The represented space is unrestricted on the $V$
part.  On the $V^\perp$ part, it contains precisely the word
coordinates in $\Lambda_t$ along the vector $P_{V^\perp}e$.
Every coordinate of $\xi_t$ has squared magnitude $n^{-t}$.
Consequently, the squared norm of the projection of the unit target
$\xi_t\otimes e$ onto $\mathcal R_t$ is
\[
  v+p_t(1-v)=p_t+(1-p_t)v.
\]
Cauchy--Schwarz applied to this projection proves
\eqref{eq:source-correlation}.
\end{proof}

\subsection{The serial recovery theorem}
\label{app:source:theorem}

\begin{theorem}[Source recovery with arbitrary cached probes]
\label{thm:source-recovery}
Fix $a\in(0,1)$, let $\rho=a^2/2$, and allow at most $M\leq H$
disclosures.  Set
\begin{equation}
  K_* =\left\lfloor\frac{\rho^2nH}{4\mathrm e}\right\rfloor,
  \qquad
  D=\left\lceil\frac{32(2K_*+1)}{a^2}\sum_{t=0}^H n^t\right\rceil.
  \label{eq:source-dimension}
\end{equation}
For the fixed product-Haar instance law in these finite dimensions,
every adaptive randomized algorithm using at most $K_*$ enlarged
oracle rounds and returning a vector $\widehat v$ with
$\|\widehat v\|\leq1$ satisfies
\begin{equation}
  \mathbb P\bigl(|\langle u_H,\widehat v\rangle|<a\bigr)
  \geq\frac{15}{16}.
  \label{eq:source-hardness}
\end{equation}
Conversely, $nH$ rounds suffice to recover $u_H$ exactly.  Thus,
for fixed $a$, the success-probability-$2/3$ terminal-source
complexity is $\Theta_a(nH)$.
\end{theorem}

\begin{proof}
Fix the private random seed and pad the algorithm to $K_*$ rounds
with zero inputs if necessary.  By
\eqref{eq:source-gaussian-projection} and the choice of $D$,
Markov's inequality gives
\[
  \mathbb P\bigl(\|P_Ve\|^2>\rho\bigr)
  \leq\frac{a^2/32}{a^2/2}=\frac1{16}.
\]
If $K_*<H$, Theorem~\ref{thm:source-counting} gives $p_H=0$.
Otherwise, using $\binom{K_*}{H}\leq(\mathrm e K_*/H)^H$
and $r_{K_*}\leq M\leq H$, it gives
\[
  p_H\leq\rho^{-M}
  \left(\frac{\mathrm e K_*}{nH}\right)^H
  \leq\rho^{-M}\left(\frac{\rho^2}{4}\right)^H
  \leq\left(\frac\rho4\right)^H
  \leq\frac\rho4.
\]
On the common event $\|P_Ve\|^2\leq\rho$,
Lemma~\ref{lem:source-correlation} therefore yields
\[
  |\langle u_H,\widehat v\rangle|^2
  \leq\frac\rho4+\rho=\frac{5a^2}{8}<a^2.
\]
Integrating over the independent private seed proves
\eqref{eq:source-hardness} for randomized algorithms under the
same fixed instance law.  In particular, a worst-case success
guarantee would also hold on this distribution, contradicting the
displayed bound.  An integer budget that succeeds must exceed
$\lfloor\rho^2nH/(4\mathrm e)\rfloor$, and hence must exceed
$\rho^2nH/(4\mathrm e)$.  This also covers $K_*=0$.

For the converse, $u_0=e$ is public.  Once $u_t$ is known, query
all $n$ branches on that vector and form
\[
  u_{t+1}=\frac1{\sqrt n}\sum_{i=1}^n A_{t,i}u_t.
\]
Applying this procedure successively to all $H$ edges uses $nH$
rounds and recovers the terminal source exactly, without any
disclosures.
\end{proof}

\begin{corollary}[Serial stages]
\label{cor:source-stages}
For $M$ stages with $m\geq1$ edges per stage, set $H=Mm$ and
allow one source disclosure at each stage endpoint.  Recovering
constant correlation with the terminal source requires
$\Omega_a(nMm)$ enlarged oracle rounds.
\end{corollary}

\begin{proof}
Apply Theorem~\ref{thm:source-recovery} with disclosure sites
$\{m,2m,\ldots,Mm\}$.  The factor $nMm$ comes from one
chronological information bound that includes all disclosures and
cached probes.
\end{proof}

\paragraph{Additional local input channels.}
The same proof accommodates at most $b$ physical input vectors per
vertex per round and at most $b_{\mathrm{out}}$ final vectors per
vertex.  These extra vectors are exposed through inverse-image
steps and receive no additional matrix actions.  Replace $2K+1$
by $bK+b_{\mathrm{out}}$ in
\eqref{eq:source-gaussian-space} and
\eqref{eq:source-dimension}.  The counting and correlation bounds
are unchanged.  Thus a fixed number of state, connector, and dual
blocks sharing a vertex rotation changes only the required finite
dimension.

\section{Quadratic relay constructions}
\label{sec:relays}

This section constructs the quadratic modules used by the lower bound.
Each module has a globally smooth component representation, although an
individual component need not be concave in the dual variable.  Maximizing
the average over the entire dual space produces a positive quadratic
interaction between two primal connectors.  We first construct a long
relay and then give a one-edge relay whose entire primal energy has a
tunable amplitude.

\subsection{The long relay}
\label{sec:long-relay}

Fix $N\ge 2$ and consecutive vertex spaces
$\mathcal H_0,\ldots,\mathcal H_{N-1}$.  For each edge, let
$A_{t,i}:\mathcal H_t\to\mathcal H_{t+1}$ be the branch maps from the
information construction, so that
\[
 A_{t,i}^{\mathsf T}A_{t,j}=\mathbf 1_{\{i=j\}}I,
 \qquad
 U_t=\frac1{\sqrt n}\sum_{i=1}^n A_{t,i},
 \qquad U_t^{\mathsf T}U_t=I.
\]
Write $P=U_{N-2}\cdots U_0$.  The primal connectors are
$a\in\mathcal H_0$ and $b\in\mathcal H_{N-1}$; the dual variable is
$w=(w_0,\ldots,w_{N-1})$, with $w_t\in\mathcal H_t$.

\subsubsection{Scalar endpoint normalization}

Let $L_N$ be the Neumann path Laplacian: its endpoint diagonal entries are
one, its interior diagonal entries are two, and its adjacent off-diagonal
entries are minus one.  Define
\begin{equation}
\begin{split}
 \lambda_N&=N^{-2},\qquad M_N=\lambda_N I+L_N,
 \qquad k_N=(M_N^{-1})_{0,N-1},\\
 c_N&=\frac{(M_N^{-1})_{0,0}}{k_N},
 \qquad d_N=\frac{2-c_N}{2}.
\end{split}
\label{eq:scalar-endpoint-normalization}
\end{equation}

\begin{lemma}[Uniform endpoint constants]
\label{lem:endpoint-constants}
For every $N\ge2$,
\[
 \frac N2<k_N\le N,
 \qquad 1\le c_N<\cosh 1<\frac85,
 \qquad \frac15<d_N\le\frac12.
\]
\end{lemma}

\begin{proof}
Put $\vartheta=2\operatorname{arsinh}(1/(2N))$.  Then
$2\cosh\vartheta-2=N^{-2}$ and $0<N\vartheta<1$.
The last column of $M_N^{-1}$ is
\[
 x_j=C\cosh((j+\tfrac12)\vartheta),\qquad
 C=\frac1{2\sinh(N\vartheta)\sinh(\vartheta/2)},
 \qquad 0\le j<N.
\]
Indeed, the root equation follows from the fictitious boundary condition
$x_{-1}=x_0$, and the interior equations follow from the hyperbolic
recurrence.  At the other endpoint, the required normalization is
\[
 \cosh((N+\tfrac12)\vartheta)
 -\cosh((N-\tfrac12)\vartheta)
 =2\sinh(N\vartheta)\sinh(\vartheta/2).
\]
Reflection symmetry makes the two diagonal endpoint Green coefficients
equal.  Since $\sinh(\vartheta/2)=1/(2N)$, we obtain
\[
 k_N=\frac{N\cosh(\vartheta/2)}{\sinh(N\vartheta)}
 >\frac N{\sinh1}>\frac N2,
 \qquad
 c_N=\frac{\cosh((N-\tfrac12)\vartheta)}
 {\cosh(\vartheta/2)}<\cosh1.
\]
Every $x_j$ is at least $x_0=k_N$.  Moreover,
$\mathbf 1^{\mathsf T}M_N=N^{-2}\mathbf 1^{\mathsf T}$ implies
$\sum_jx_j=N^2$, whence $Nk_N\le N^2$.  Finally,
$c_N\ge1$ and $\cosh1<8/5$ give the stated bounds on $d_N$.
\end{proof}

Set
\begin{equation}
 \beta=n^{-1/2},\qquad r_N=\sqrt{\beta/k_N}.
 \label{eq:long-relay-coefficients}
\end{equation}
For each component $i$, define
\begin{equation}
\begin{split}
 \mathcal H_i(a,b;w)
 ={}&d_N(\|a\|^2+\|b\|^2)
       +r_N\bigl(\langle a,w_0\rangle
                   -\langle b,w_{N-1}\rangle\bigr)\\
 &-\frac\beta2\left[
      \lambda_N\sum_{t=0}^{N-1}\|w_t\|^2
      +\sum_{t=0}^{N-2}(\|w_t\|^2+\|w_{t+1}\|^2)
   \right]
   +\sum_{t=0}^{N-2}\langle w_{t+1},A_{t,i}w_t\rangle.
\end{split}
\label{eq:long-component}
\end{equation}
Its average is exactly
\begin{equation}
\begin{split}
 \mathcal H(a,b;w)
 ={}&d_N(\|a\|^2+\|b\|^2)
       +r_N\bigl(\langle a,w_0\rangle
                   -\langle b,w_{N-1}\rangle\bigr)\\
 &-\frac\beta2\left[
     \lambda_N\sum_{t=0}^{N-1}\|w_t\|^2
     +\sum_{t=0}^{N-2}\|w_{t+1}-U_tw_t\|^2
   \right].
\end{split}
\label{eq:long-average}
\end{equation}

\begin{theorem}[The long relay]
\label{thm:long-relay}
The components in \eqref{eq:long-component} have the following properties.
\begin{enumerate}
\item Each component has a globally $8$-Lipschitz joint gradient,
independently of $n$, $N$, and the dimensions.  Their average is
$\beta/N^2$-strongly concave in $w$.
\item There is a self-adjoint operator $\mathcal D_N$ on
$\mathcal H_{N-1}$ such that
\begin{equation}
\begin{split}
 Q(a,b):=\max_w\mathcal H(a,b;w)
 &=\|a\|^2-\langle b,Pa\rangle
       +\langle b,\mathcal D_Nb\rangle,\\
 d_N I\preceq\mathcal D_N&\preceq I,
 \qquad \mathcal D_NP=P.
\end{split}
\label{eq:long-endpoint-energy}
\end{equation}
Consequently, for $(a,b)\ne(0,0)$,
\begin{equation}
 \frac15(\|a\|^2+\|b\|^2)
 <Q(a,b)\le\frac32(\|a\|^2+\|b\|^2).
 \label{eq:long-coercivity}
\end{equation}
In particular, $Q$ has a dimension-independent positive Hessian lower
bound.  On the aligned source line
$a=a_0u_0$, $b=b_0u_{N-1}$, its value is
$a_0^2-a_0b_0+b_0^2$.
\item The unique unconstrained dual maximizer satisfies
\begin{equation}
 \|w^*(a,b)\|
 \le\sqrt3\,Nn^{1/4}\sqrt{\|a\|^2+\|b\|^2}.
 \label{eq:long-dual-radius}
\end{equation}
\item One value-and-gradient query to a selected component is simulated
exactly by one forward and one transpose action of its selected branch
map on each internal edge, followed by public arithmetic.  The component
formula contains no hidden source vector.
\end{enumerate}
\end{theorem}

\begin{proof}
Let $\mathcal M$ denote the operator in the square bracket of
\eqref{eq:long-average}.  We have $\mathcal M\succeq N^{-2}I$.
Equality is attained on a nonzero tuple satisfying
$w_{t+1}=U_tw_t$, so the smallest eigenvalue is exactly $N^{-2}$.
This proves the average strong-concavity assertion.
For a component Hessian, the diagonal blocks have norm at most three,
the symmetric off-diagonal dual part has norm at most two, and the
endpoint primal--dual coupling has norm at most $2r_N\le2$.
The last inequality follows from $k_N>N/2\ge1$ and $\beta\le1$.
The sum is at most seven and is therefore bounded by eight.

To compute the endpoint blocks of $\mathcal M^{-1}$, we must also account
for directions orthogonal to the transmitted source.  Define the
innovation spaces
\[
 \mathcal K_0=\mathcal H_0,
 \qquad \mathcal K_j=\ker U_{j-1}^{\mathsf T}\quad(j\ge1).
\]
Because every $U_t$ is an isometry,
\[
 \mathcal H_t
 =\bigoplus_{j=0}^t U_{t-1}\cdots U_j\mathcal K_j,
\]
where the product is the identity when $j=t$.
This orthogonal decomposition separates $\mathcal M$ into scalar suffix
chains.  The $j=0$ chain is $M_N$.  For $j>0$, it is the trailing
principal submatrix of $M_N$: the incoming edge adds a unit anchoring term
at the first vertex of that suffix.  It follows that
\begin{equation}
 (\mathcal M^{-1})_{0,0}=c_Nk_NI,
 \qquad
 (\mathcal M^{-1})_{N-1,0}=k_NP.
 \label{eq:long-root-green}
\end{equation}
Writing $G=(\mathcal M^{-1})_{N-1,N-1}$, we also have
\begin{equation}
 0\prec G\preceq c_Nk_NI,
 \qquad GP=c_Nk_NP.
 \label{eq:long-terminal-green}
\end{equation}
For the upper bound, the inverse of a positive principal submatrix is
bounded above by the corresponding block of the inverse of the full
matrix.  This is the Schur-complement inequality.  Apply it to each
scalar suffix chain and then use the orthogonal decomposition above.
On the root-generated chain, equality holds, which gives the identity
for $GP$.

Let $E(a,b)$ insert $a$ at layer zero and $-b$ at layer $N-1$.
Completing the square yields
\[
 Q(a,b)=d_N(\|a\|^2+\|b\|^2)
       +\frac1{2k_N}
        \langle E(a,b),\mathcal M^{-1}E(a,b)\rangle.
\]
Equations \eqref{eq:long-root-green}--\eqref{eq:long-terminal-green}
give \eqref{eq:long-endpoint-energy} with
$\mathcal D_N=d_NI+G/(2k_N)$, since $d_N+c_N/2=1$.
To verify coercivity, decompose $b=Pc+b_\perp$, where
$P^{\mathsf T}b_\perp=0$.  Self-adjointness and
$\mathcal D_NP=P$ give
\[
 Q(a,b)=\|a\|^2-\langle a,c\rangle+\|c\|^2
          +\langle b_\perp,\mathcal D_Nb_\perp\rangle.
\]
The first three terms lie between one half and three halves of
$\|a\|^2+\|c\|^2$.  Combining this with
$d_N>1/5$ and $\mathcal D_N\preceq I$ proves
\eqref{eq:long-coercivity}.  The source-line formula follows from
$Pu_0=u_{N-1}$.

The dual maximizer is
\[
 w^*(a,b)=\frac1{\sqrt{\beta k_N}}
                  \mathcal M^{-1}E(a,b).
\]
Since $\mathcal M\succeq N^{-2}I$, spectral calculus implies
$\mathcal M^{-2}\preceq N^2\mathcal M^{-1}$.
The endpoint compression of $\mathcal M^{-1}$ has norm at most
$k_N(c_N+1)$ by
\eqref{eq:long-root-green}--\eqref{eq:long-terminal-green}.  Hence
\begin{align*}
 \|w^*(a,b)\|^2
 &\le\frac{N^2}{\beta k_N}
       \langle E(a,b),\mathcal M^{-1}E(a,b)\rangle\\
 &\le\frac{(c_N+1)N^2}{\beta}
                  (\|a\|^2+\|b\|^2)\\
 &<3N^2n^{1/2}(\|a\|^2+\|b\|^2),
\end{align*}
which proves \eqref{eq:long-dual-radius}.

Finally, every hidden matrix in \eqref{eq:long-component} occurs in a
local bilinear term.  One forward and one transpose action per edge
provide its value and both gradients.  All remaining terms are public.
\end{proof}

\subsection{A one-edge relay with adjustable amplitude}
\label{sec:short-relay}

The long relay has condition number of order $\sqrt n\,N^2$.
For smaller condition numbers, we retain a single branching edge and
adjust the scalar normalization of the eliminated primal energy.
Let $A_i:\mathcal H_0\to\mathcal H_1$ be one layer of branch maps, and set
\[
 \beta=n^{-1/2},\qquad
 U=\frac1{\sqrt n}\sum_{i=1}^nA_i,
 \qquad U^{\mathsf T}U=I.
\]
For any $\mu_0\in(0,1]$, define
\begin{equation}
 t=\mu_0+\beta,
 \qquad\theta=\frac\beta t,
 \qquad r^2=t-\frac{\beta^2}{t},
 \qquad d=\theta-\frac12.
 \label{eq:short-relay-coefficients}
\end{equation}
For primal connectors $a,b$ and dual variables $w_0,w_1$, use
\begin{equation}
\begin{split}
 \mathcal J_i(a,b;w_0,w_1)
 ={}&d(\|a\|^2+\|b\|^2)
       +r\bigl(\langle a,w_0\rangle-\langle b,w_1\rangle\bigr)\\
 &-\frac t2(\|w_0\|^2+\|w_1\|^2)
       +\langle w_1,A_iw_0\rangle.
\end{split}
\label{eq:short-component}
\end{equation}
The coefficient $d$ may be negative, consistently with the absence of a
primal convexity assumption.  After maximizing the average, however,
the value divided by $\theta$ is uniformly positive definite.

\begin{lemma}[The one-edge relay]
\label{lem:short-relay}
Every component in \eqref{eq:short-component} has a globally
$8$-Lipschitz joint gradient.  Its average is $\mu_0$-strongly concave
in the entire dual space, and its exact value is
\begin{equation}
\begin{split}
 \max_{w_0,w_1}\frac1n\sum_{i=1}^n
             \mathcal J_i(a,b;w_0,w_1)
 &=\theta\bigl[\|a\|^2-\langle b,Ua\rangle
                         +\langle b,D_\theta b\rangle\bigr],\\
 D_\theta&=I-\frac\theta2(I-UU^{\mathsf T}).
\end{split}
\label{eq:short-endpoint-energy}
\end{equation}
In particular, $\tfrac12I\preceq D_\theta\preceq I$ and
$D_\theta U=U$.  At $a=b=0$, the unique dual maximizer is zero.
A component query is simulated by one forward and one transpose action
of its selected $A_i$, together with public arithmetic.
\end{lemma}

\begin{proof}
The negative dual Hessian of the average is
\[
 H=\begin{pmatrix}
 tI&-\beta U^{\mathsf T}\\
 -\beta U&tI
 \end{pmatrix}.
\]
On $\mathcal H_0\oplus\operatorname{range}(U)$ its eigenvalues are
$t-\beta=\mu_0$ and $t+\beta$.  On the terminal complement
$\operatorname{range}(U)^\perp$, it acts as $tI$.
Thus $H$ is positive definite with minimum eigenvalue $\mu_0$.
Its inverse blocks are
\begin{align*}
 (H^{-1})_{00}&=\frac t{t^2-\beta^2}I,
 & (H^{-1})_{10}&=\frac\beta{t^2-\beta^2}U,\\
 (H^{-1})_{11}&=\frac1tI
          +\frac{\beta^2}{t(t^2-\beta^2)}UU^{\mathsf T}.
\end{align*}
Completing the square adds
$\tfrac{r^2}{2}\langle(a,-b),H^{-1}(a,-b)\rangle$
to the raw primal quadratic.  Using
$r^2=(t^2-\beta^2)/t$ and $d=\beta/t-1/2$, the root coefficient becomes
$d+1/2=\theta$, and the cross coefficient becomes
$-r^2\beta/(t^2-\beta^2)=-\theta$.
The terminal coefficient is
\begin{align*}
 dI+\frac{r^2}{2t}I
       +\frac{r^2\beta^2}{2t(t^2-\beta^2)}UU^{\mathsf T}
 &=\theta I-\frac{\theta^2}{2}(I-UU^{\mathsf T})\\
 &=\theta D_\theta.
\end{align*}
This proves \eqref{eq:short-endpoint-energy}.  The assertions about
$D_\theta$ follow because $UU^{\mathsf T}$ is an orthogonal projector
and $0<\theta<1$.

For component smoothness, observe that
$|2d|\le1$, $t\le2$, $r\le\sqrt2$, and $\|A_i\|=1$.
The diagonal Hessian part has norm at most two, the dual off-diagonal
part has norm at most one, and the connector--dual coupling has norm
at most $r$.  Their sum is less than eight.  This estimate holds
globally.  At zero connectors, the average equals
$-\langle w,Hw\rangle/2$, whose unique maximizer is zero.
The oracle statement follows directly from the single hidden bilinear
term in \eqref{eq:short-component}.
\end{proof}

The scalar correction $d=\theta-1/2$ is the key to this normalization.
Without it, elimination produces an endpoint self-energy of order one
and a transmitted cross term of order $\theta$.
The correction makes the entire value, including the terminal
orthogonal directions, have amplitude $\theta$.
The projector $UU^{\mathsf T}$ appears only in the eliminated value;
it is absent from the component formula and therefore is not exposed
by a component query.

\section{Global construction and regularity}
\label{sec:construction}

We next assemble the relay modules into a serial instance.
All nonlinear terms will depend only on primal variables.  This
separation permits arbitrary dual query norms while retaining global
component smoothness.

\subsection{Scalar profiles and radial caps}
\label{sec:profiles-caps}

Define the smooth monotone step
\[
 \omega(t)=
 \begin{cases}
 0,&t\le0,\\
 \displaystyle\frac{e^{-1/t}}{e^{-1/t}+e^{-1/(1-t)}},&0<t<1,\\
 1,&t\ge1.
 \end{cases}
\]
It is flat at both endpoints, satisfies
$\omega(t)+\omega(1-t)=1$, and obeys
\begin{equation}
 0\le\omega'\le8,\qquad |\omega''|\le128.
 \label{eq:omega-bounds}
\end{equation}
For an explicit verification, take $0<t\le1/2$ and put $u=1/t\ge2$.
The logistic derivative gives
\[
 \omega'
 =\omega(1-\omega)\bigl[t^{-2}+(1-t)^{-2}\bigr],
 \qquad \omega(1-\omega)\le e^{2-u}.
\]
Thus $\omega'\le e^{2-u}(u^2+4)\le8$.
Differentiating once more gives
\[
 |\omega''|
 \le e^{2-u}(u^4+2u^3+8u^2+32)<128.
\]
The last inequality follows by bounding separately the four monomials
times $e^{-u}$ on $[2,\infty)$.
Symmetry treats $1/2\le t<1$.

Let $p(t)=0$ for $t\le0$ and
$p(t)=\int_0^t\omega(v)\,dv$ for $t>0$.  Symmetry implies
$p(t)=t-1/2$ for $t\ge1$.  Set
\begin{equation}
 q(t)=\omega((5t-1)/4),\qquad
 e_s(t)=p(t+3)-p(t-2)-\frac52.
 \label{eq:scalar-profiles}
\end{equation}
Then $q=0$ on $(-\infty,1/5]$, $q=1$ on $[1,\infty)$,
$0\le q'\le10$, and $|q''|\le200$.
Also, $e_s(t)=t$ on $[-2,2]$, $|e_s|\le5/2$,
$0\le e_s'\le1$, $|e_s''|\le16$, and $e_s'=0$ outside $[-3,3]$.

For $C>0$, define
\[
 f_C(r)=
 \begin{cases}
 r,&0\le r\le C,\\
 \displaystyle C+C\int_0^{(r-C)/C}(1-\omega(v))\,dv,
     &C<r<2C,\\
 3C/2,&r\ge2C.
 \end{cases}
\]
On any Euclidean space, let
\begin{equation}
 \rho_C(x)=
 \begin{cases}
 f_C(\|x\|)x/\|x\|,&x\ne0,\\
 0,&x=0.
 \end{cases}
 \label{eq:radial-cap}
\end{equation}
This smooth map equals the identity when $\|x\|\le C$ and is always a
scalar multiple of $x$.  Its bounds are independent of the dimension:
\begin{equation}
 \|\rho_C(x)\|\le\min\{\|x\|,2C\},\qquad
 \|D\rho_C(x)\|\le1,\qquad
 \|D^2\rho_C(x)\|\le18/C.
 \label{eq:radial-cap-bounds}
\end{equation}
Indeed, $0\le f_C'\le1$, $f_C(r)\le r$, $f_C(r)\ge C$ for $r\ge C$,
and $|f_C''|\le8/C$.  With $g(r)=f_C(r)/r$, differentiation gives
\[
 \|D^2(g(\|x\|)x)\|
 \le3|g'(r)|+r|g''(r)|
 \le|f_C''(r)|+\frac{10}r\qquad(r>C),
\]
which implies the last bound in \eqref{eq:radial-cap-bounds}.
Flatness of $\omega$ makes the pieces glue smoothly.
In particular, for any unit vector $u$, the score
$x\mapsto\langle u,\rho_C(x)\rangle$ is $1$-Lipschitz.

\subsection{A common outer construction for both relays}
\label{sec:outer-construction}

Let $H=Mm$, where $M,m\ge1$, and place a source site at each depth
$jm$.  Write $v_j=u_{jm}$ for $1\le j\le M$ and $v_0=e$.
Stage $j$ has connectors
$a_j\in\mathcal H_{(j-1)m}$ and $b_j\in\mathcal H_{jm}$,
and a state $S_j\in\mathcal H_{jm}$.
Let $P_j$ be the product of the $m$ average isometries in that stage.
The only properties required of its eliminated quadratic are
\begin{equation}
\begin{split}
 Q_j(a,b)&=\|a\|^2-\langle b,P_ja\rangle
                        +\langle b,D_jb\rangle,\\
 \frac15I\preceq D_j&\preceq I,\qquad D_jP_j=P_j,
 \qquad P_jv_{j-1}=v_j.
\end{split}
\label{eq:endpoint-geometry}
\end{equation}
The module has value $\theta Q_j$.
For the long relay, $m=N-1$ and $\theta=1$.
For the one-edge relay, $m=1$ and $\theta$ is given by
\eqref{eq:short-relay-coefficients}.
All stages in one instance use the same amplitude.

\subsubsection{State penalties and couplings}

Fix the numerical constants
\begin{equation}
\begin{gathered}
 C_\nu=21,\qquad \sigma=42,\qquad \gamma=4300,
 \qquad K_s=172,\qquad R_s=10,\\
 \delta_g=\frac1{430000},\qquad c_R=32416,
 \qquad \tau_* =\frac1{20}.
\end{gathered}
\label{eq:outer-constants}
\end{equation}
Put $E_\nu=\rho_{21}$ and $Z=\rho_{10}$.  Define
\begin{equation}
\begin{split}
 t_j&=\langle v_j,Z(S_j)\rangle,\qquad t_0=1,\\
 h(t)&=t\,\omega((t-\delta_g)/\delta_g),\\
 V(t)&=-\frac{c_R+1}{10}
                 \bigl[p(10t-1)-p(10t-10)\bigr],\\
 P_j^{\mathrm{st}}(S_j)
 &=\frac\gamma2\|S_j\|^2-\frac\gamma2h(t_j)^2+V(t_j).
\end{split}
\label{eq:outer-state}
\end{equation}
The notation $P_j^{\mathrm{st}}$ denotes the state penalty and is
distinct from the isometry $P_j$ in \eqref{eq:endpoint-geometry}.
The three couplings are
\begin{align}
 G_{\mathrm{en},j}
 &=-4q(t_{j-1})(1-q(t_j))
                      \langle v_{j-1},E_\nu(a_j)\rangle,
 \label{eq:outer-entrance}\\
 G_{\mathrm{ex},j}
 &=-q(t_{j-1})(1-q(t_j))
                      \langle Z(S_j),E_\nu(b_j)\rangle,
 \label{eq:outer-exit}\\
 G_{\mathrm{st},j}
 &=K_sq(t_j)(1-e_s(t_{j-1}))(1-q(t_{j-1})).
 \label{eq:outer-coupling}
\end{align}
The root data $t_0=1$ and $v_0=e$ are public.
In particular, the exit term uses an ordinary vector inner product;
it does not apply a hidden rank-one projection to the terminal
connector.
For $\nu=((a_j,b_j))_{j=1}^M$, let
\begin{equation}
\begin{split}
 R(S,\nu)&=\sum_{j=1}^M
   \bigl(G_{\mathrm{en},j}+G_{\mathrm{ex},j}
                 +G_{\mathrm{st},j}+P_j^{\mathrm{st}}\bigr),\\
 F_Q(S,\nu)&=R(S,\nu)+\sum_{j=1}^M Q_j(a_j,b_j).
\end{split}
\label{eq:outer-objective}
\end{equation}
The raw components are
\begin{equation}
 f_i^0(S,\nu,w)
   =\sum_{j=1}^M\mathcal K_i^{(j)}(a_j,b_j;w^{(j)})
       +\theta R(S,\nu),
 \label{eq:assembled}
\end{equation}
where $\mathcal K_i^{(j)}$ is either the long module
$\mathcal H_i$ or the short module $\mathcal J_i$.
Different stages have distinct dual tuples, including when their
endpoints occupy the same vertex space.
Maximization over the full dual space gives the exact identity
\begin{equation}
 \Phi_0(S,\nu)
 :=\max_{w\in\mathbb R^{d_y}}\frac1n\sum_{i=1}^n
                    f_i^0(S,\nu,w)
 =\theta F_Q(S,\nu).
 \label{eq:value-identity}
\end{equation}

\begin{lemma}[Global admissibility]
\label{lem:global-admissibility}
Every component in \eqref{eq:assembled} is jointly
$L_*=10^8$-smooth on its entire Euclidean domain.
The average dual strong-concavity constant is $\beta/N^2$ for the long
family and $\mu_0$ for the short family.
Moreover, $y^*(0)=0$, $\Phi_0(0)=0$, and
\begin{equation}
 \Phi_0(0)-\inf\Phi_0\le\theta C_\Delta M,
 \qquad C_\Delta=31000.
 \label{eq:normalized-gap}
\end{equation}
\end{lemma}

\begin{proof}
The relay modules are globally eight-smooth and block diagonal across
stages.  We give explicit derivative bounds for the remaining terms;
their bounded incidence will make the final estimate independent of
the number of stages.

On the score range $[-20,20]$,
$\|Dt_j\|\le1$ and $\|D^2t_j\|\le2$.
Thus the composed profile $q(t_j)$ has gradient and Hessian bounds
$10$ and $220$, respectively, while $e_s(t_j)$ has bounds $1$ and $18$.
The connector cap has image bound $42$, derivative bound one, and
second-derivative bound one.  The state cap has image bound $20$ and
second-derivative bound two.
For the remaining scalar terms,
\[
 |(h^2)'|\le680,\qquad |(h^2)''|\le1666,
 \qquad |V'|\le32417,
 \qquad |V''|\le160\cdot32417.
\]
To see why the narrow transition of $h$ is harmless, within that
transition we have
$|h|\le2\delta_g$, $|h'|\le17$, and
$|h''|\le272/\delta_g$.
Consequently,
$|(h^2)''|\le2\cdot17^2+4\delta_g(272/\delta_g)=1666$.
Outside the transition, $h$ is either zero or the identity.
After composition with the score, these inequalities bound the Hessian
of every state penalty by $12\cdot10^6$.

The product rule gives the following local Hessian bounds:
\[
 \|\nabla^2G_{\mathrm{en},j}\|\le2\cdot10^5,
 \qquad
 \|\nabla^2G_{\mathrm{ex},j}\|\le6\cdot10^5,
 \qquad
 \|\nabla^2G_{\mathrm{st},j}\|\le4\cdot10^5.
\]
For example, the exit inner product has absolute value at most $840$,
gradient norm at most $62$, and Hessian norm at most $106$.
Its gate product has derivative norm at most $20$ and Hessian norm at
most $640$.  Therefore its Hessian is bounded by
$840\cdot640+2\cdot20\cdot62+106<6\cdot10^5$.
The entrance bound follows in the same way from the capped linear
functional, with bounds $42$, $1$, and $1$.
For the state coupling, the factor
$(1-e_s(t_{j-1}))(1-q(t_{j-1}))$ has value, gradient, and Hessian bounds
$7/2$, $36$, and $808$.  Hence a bound is
$172[(7/2)\cdot220+808+2\cdot10\cdot36]<4\cdot10^5$.

Each state is incident only to its own penalty and the two neighboring
stages.  To formalize the resulting operator bound, let $H_r$ be local
Hessians supported on block sets $I_r$.  For any vector $z$,
\[
 \left|\left\langle z,\sum_r H_rz\right\rangle\right|
 \le\sum_r\|H_r\|\,\|z_{I_r}\|^2
 \le\left(\max_b\sum_{r:\,b\in I_r}\|H_r\|\right)\|z\|^2.
\]
The maximum incident sum is below $4\cdot10^7$, so
$\|\nabla^2R\|<4\cdot10^7$ independently of $M$, $n$, $N$, and
the dimension.  Since $0<\theta\le1$, adding the relay modules is
covered by $L_*=10^8$.
All added terms are purely primal, so average dual strong concavity
follows directly from the corresponding relay lemma.

For the gap bound, write $b=P_jc+b_\perp$ with
$P_j^{\mathsf T}b_\perp=0$.  From
\eqref{eq:endpoint-geometry},
\[
 Q_j(a,b)
 =\|a\|^2-\langle a,c\rangle+\|c\|^2
          +\langle b_\perp,D_jb_\perp\rangle
 \ge\frac15(\|a\|^2+\|b\|^2).
\]
Thus $\sum_jQ_j$ is nonnegative and its Hessian is at least $2I/5$.
Also, $|h(t_j)|\le|t_j|\le\|S_j\|$, so the two quadratic terms
in $P_j^{\mathrm{st}}$ have nonnegative sum.
The state coupling is nonnegative: whenever $q(t)<1$,
we have $e_s(t)\le1$.
The absolute sizes of the remaining negative contributions per stage
are bounded by
\[
 \frac9{10}(32417)+4\sigma+2R_s\sigma<31000.
\]
This proves \eqref{eq:normalized-gap} from the eliminated value
\eqref{eq:value-identity}.  The argument does not require the raw
primal quadratic in the short module to be coercive.
At the primal origin, all outer terms and connectors vanish.
The average dual objective is then a strictly concave quadratic with
unique maximizer zero, proving $y^*(0)=0$ and $\Phi_0(0)=0$.
Finally, the explicit eliminated formula establishes smoothness of
the primal value; no uniform bound on its smoothness constant is
needed here.
\end{proof}

\section{Stationarity and exact oracle simulation}
\label{app:stationarity}

This section completes the interface between the analytic construction and
the information problem. First, a small gradient of the eliminated primal
objective forces the terminal score to be high. Second, the original
component values and full gradients can be simulated exactly by the
source-release oracle until that recovery occurs.

\begin{samepage}
\subsection{A direct primal-gradient certificate}

\begin{theorem}[Direct primal-gradient certificate]
\label{thm:gradient-certificate}
For every choice of isometries and self-adjoint operators satisfying
\eqref{eq:endpoint-geometry}, the function $F_Q$ in
\eqref{eq:outer-objective} satisfies
\begin{equation}
 \|\nabla F_Q(S,\nu)\|\le\tau_*
 \quad\Longrightarrow\quad t_M\ge1.
 \label{eq:gradient-certificate}
\end{equation}
Consequently, $\|\nabla\Phi_0(S,\nu)\|\le\theta\tau_*$ forces terminal
source recovery. The implication holds for arbitrary primal points, with
no restriction to source-aligned directions or bounded domains.
\end{theorem}
\end{samepage}

\begin{proof}
Let $G$ be the sum of the entrance, exit, and state couplings in
\eqref{eq:outer-entrance}--\eqref{eq:outer-coupling}, excluding the state
penalties. We use the constants in \eqref{eq:outer-constants} throughout.
Assume $\|\nabla F_Q\|\le1/20$.

\paragraph{Step 1: isolate the scalar and vector derivatives.}
Separate the direct occurrence of $Z(S_j)$ in the exit inner product from
its occurrences through the score $t_j$. Then
\begin{equation}
\begin{split}
 \nabla_{S_j}F_Q
 ={}&\gamma S_j+B_j\nabla t_j\\
 &-q(t_{j-1})(1-q(t_j))\,DZ(S_j)^{\mathsf T}E_\nu(b_j),
\end{split}
\label{eq:state-gradient}
\end{equation}
where
\[
 B_j=-\gamma h(t_j)h'(t_j)+V'(t_j)+C_j,
\]
and $C_j$ is the score derivative of $G$, holding its direct vector
occurrences fixed. For clarity, put
\[
 \ell_j=\langle v_{j-1},E_\nu(a_j)\rangle,
 \qquad b_j^{\mathrm{ex}}=\langle Z(S_j),E_\nu(b_j)\rangle,
 \qquad r_0(t)=(1-e_s(t))(1-q(t)).
\]
For $j=M$, terms with index $M+1$ below are absent. Direct differentiation
gives
\begin{align*}
 C_j={}&4q'(t_j)\bigl[q(t_{j-1})\ell_j
                   -(1-q(t_{j+1}))\ell_{j+1}\bigr]\\
 &+q'(t_j)\bigl[q(t_{j-1})b_j^{\mathrm{ex}}
                   -(1-q(t_{j+1}))b_{j+1}^{\mathrm{ex}}\bigr]\\
 &+K_s\bigl[q'(t_j)r_0(t_{j-1})+q(t_{j+1})r_0'(t_j)\bigr].
\end{align*}
Using $|\ell_j|\le\sigma$, $|b_j^{\mathrm{ex}}|\le2R_s\sigma$,
$|1-e_s|\le7/2$, $0\le e_s'\le1$, and $0\le q'\le10$, we obtain
\begin{align}
 |C_j|&\le K_s(1+7\cdot10)+8\cdot10\sigma
                     +4R_s\cdot10\sigma=32372,\label{eq:C-bound}\\
 |t_jC_j|&\le K_s(3+7\cdot10)+8\cdot10\sigma
                     +4R_s\cdot10\sigma=32716
                     \quad(t_j\le1).\label{eq:tC-bound}
\end{align}
For the second inequality, $q'$ is supported on $[1/5,1]$ and
$|te_s'(t)|\le3$. Also, for $t\le1$,
$|th(t)h'(t)|\le17$ and $|tV'(t)|\le32417$. Hence
\begin{equation}
 |t_jB_j|\le4300\cdot17+32417+32716=138233
 \qquad(t_j\le1).
 \label{eq:tB-bound}
\end{equation}

\paragraph{Step 2: localize unfinished states.}
Suppose $t_j\le1$ and $\|S_j\|=r\ge R_s=10$. The radial cap has
$Z(S_j)=f_{10}(r)S_j/r$, with $f_{10}(r)\ge10$ and
$0\le f_{10}'(r)\le1$. Along the radial direction,
\[
 \frac{dt_j}{dr}=t_j\frac{f_{10}'(r)}{f_{10}(r)}.
\]
Equations~\eqref{eq:state-gradient} and \eqref{eq:tB-bound} therefore
bound the radial derivative below by
\[
 4300r-\frac{138233}{10}-42\ge29134.7>1.
\]
This contradicts the assumed gradient bound. Thus every state with
$t_j\le1$ lies in the identity region of $Z$, and
$t_j=\langle v_j,S_j\rangle$. This step uses only capped connector bounds,
so it does not presuppose bounded connectors.

\paragraph{Step 3: exclude intermediate scores.}
If $t_j\le-1/10$, then $q=q'=h=V'=0$. The derivative of the outgoing
state coupling is nonpositive, and the direct exit contributes at most
$42$. The derivative along $v_j$ is consequently at most
$4300t_j+42\le-388$, a contradiction.

If $1/5<t_j<1$, then $h(t_j)=t_j$, and the derivatives of the quadratic
state terms cancel. On this interval $V'(t_j)=-32417$. By
\eqref{eq:C-bound}, the derivative along $v_j$ is at most
\[
 -32417+32372+42=-3,
\]
again a contradiction. Each state is therefore either
\emph{low}, with $-1/10<t_j\le1/5$, or \emph{high}, with $t_j\ge1$.

For a low score, define its scalar penalty derivative by
\[
 D_0(t)=\gamma t-\gamma h(t)h'(t)+V'(t).
\]
It satisfies
\begin{equation}
 D_0(t)\le2\gamma\delta_g=\frac1{50}.
 \label{eq:low-penalty}
\end{equation}
Indeed, for $t\le\delta_g$ we have $h=0$ and $V'\le0$; for
$\delta_g<t<2\delta_g$, $hh'\ge0$ and $t<2\delta_g$; and for
$t\ge2\delta_g$ the quadratic derivatives cancel and $V'\le0$.

\paragraph{Step 4: identify the single unfinished frontier.}
A low state cannot precede a high state. In that configuration the outgoing
state-coupling derivative is $-K_s=-172$; the remaining score derivatives
vanish, and the direct exit contributes at most $42$. The directional
derivative is at most
\[
 \frac1{50}+42-172<-1.
\]
Suppose for contradiction that $t_M\le1/5$. Since the public root has
$t_0=1$, the high states form an initial segment followed by low states.
There is a unique high-to-low frontier, say stage $j$. Only this stage has
active connector forcing.

Projecting its state gradient onto $v_j^\perp$ gives
\[
 P_{v_j^\perp}\nabla_{S_j}F_Q
 =4300P_{v_j^\perp}S_j-P_{v_j^\perp}E_\nu(b_j).
\]
It follows that
\[
 \|P_{v_j^\perp}S_j\|
 \le\frac{42+1/20}{4300}<\frac1{100},
 \qquad \|S_j\|<\frac14.
\]
The entrance and exit connector forcing norms are thus at most $4$ and
$1/4$, respectively.

\paragraph{Step 5: remove the connector caps and obtain a contradiction.}
The quadratic $\sum_jQ_j$ has Hessian at least $2I/5$ and gradient zero at
the origin. Its strong monotonicity yields
\[
 \|\nu\|\le\frac52\bigl(\|\nabla_\nu F_Q\|+
                               \|\nabla_\nu G\|\bigr)
 \le\frac52\left(\frac1{20}+4+\frac14\right)
 =\frac{43}{4}<21.
\]
All connector caps are therefore identities. Write
$a_{\mathrm{sc}}=\langle v_{j-1},a_j\rangle$,
$b_{\mathrm{sc}}=\langle v_j,b_j\rangle$, and $s=t_j$.
Since $P_jv_{j-1}=v_j$ and $D_jv_j=v_j$, the projected connector gradients
satisfy
\[
 g_a=2a_{\mathrm{sc}}-b_{\mathrm{sc}}-4,
 \qquad g_b=-a_{\mathrm{sc}}+2b_{\mathrm{sc}}-s,
 \qquad |g_a|,|g_b|\le\frac1{20}.
\]
Solving these two equations gives
\[
 b_{\mathrm{sc}}=\frac{4+2s+g_a+2g_b}{3}
 >\frac{4-1/5-3/20}{3}=\frac{73}{60}.
\]
The state derivative along $v_j$ is $D_0(s)-b_{\mathrm{sc}}$, which by
\eqref{eq:low-penalty} is less than $1/50-73/60<-1$. This contradiction
excludes a low terminal state. The conclusion for $\Phi_0$ follows from
the exact identity \eqref{eq:value-identity}.
\end{proof}

\begin{remark}[Existence of stationary points]
Taking $S_j=2v_j$ at every stage and all connectors zero gives a stationary
point of $F_Q$. The average dual maximizer is also zero there. Thus the
certificate does not rely on the absence of stationary points.
\end{remark}

\subsection{Exact simulation of unrestricted component queries}

\begin{theorem}[Distributional direct-gradient hardness]
\label{thm:distributional-hardness}
Fix $M,m\ge1$ and either assembled family \eqref{eq:assembled}. There is
a universal $c_{\mathrm{inf}}>0$ and a fixed distribution of instances in
finite dimensions such that every unrestricted randomized algorithm using
at most $\lfloor c_{\mathrm{inf}}nMm\rfloor$ component queries satisfies
\begin{equation}
 \PP\bigl(\|\nabla\Phi_0(\widehat x)\|\le\theta\tau_*\bigr)
 \le\frac1{16}.
 \label{eq:distributional-hardness}
\end{equation}
Queries may have arbitrary norms and return exact values and full gradients.
The output need not have been queried, and the instance law is independent
of the private seed.
\end{theorem}

\begin{proof}
Use the source-release oracle of Appendix~\ref{app:source}, with source
sites $jm$ and
\begin{equation}
 a_* =\frac{\delta_g}{2R_s}=\frac1{8600000},\qquad
 \rho=\frac{a_*^2}{2},\qquad
 c_{\mathrm{inf}}=\frac{\rho^2}{4\mathrm e},\qquad H=Mm.
 \label{eq:information-constants}
\end{equation}
Put $K_* =\lfloor c_{\mathrm{inf}}nH\rfloor$ and choose
\begin{equation}
 D=\left\lceil\frac{32(8K_*+8)}{a_*^2}
                    \sum_{t=0}^{H}n^t\right\rceil.
 \label{eq:optimization-dimension}
\end{equation}
At a shared endpoint depth, at most five physical vectors enter a query:
a state, two adjacent connectors, and two dual endpoints. Every radial cap
is a scalar multiple of its input. Eight input slots per vertex per round,
and eight final-output slots, therefore suffice for the additional-channel
extension of the source theorem. These slots increase the analysis
dimension but do not add matrix actions.

Fix the private seed. A deterministic algorithm is run on simulated
numerical component replies. Process all query inputs before matrix replies.
After the matrix round, release eligible sources \emph{exhaustively}, before
evaluating the primal part of the simulated reply. This is an allowed
release schedule: at its end every unreleased site has $p_{jm}<\rho$, and
whenever a source is released all preceding source sites are released as
well. The formal-coordinate exposure remains analysis-side information.

On the common Gaussian event in the source proof,
$\|P_Ve\|^2\le\rho$. For every unreleased source, the correlation estimate
applied to the bounded score input gives
\[
 \frac{|t_j|}{2R_s}
 =\left|\left\langle v_j,\frac{Z(S_j)}{2R_s}\right\rangle\right|
 <\sqrt{2\rho}=a_*,
 \qquad |t_j|<\delta_g.
\]
Consequently $q,q',h,h',V,V'$ vanish at that score. Complementary gates
such as $1-q$ reduce to known constants. We check each source-dependent
term in the raw component formula.
\begin{itemize}
 \item The source-dependent part of the state penalty has zero value and
 derivative; its isotropic quadratic part is public.
 \item An entrance vanishes with its full first derivative when its
 predecessor is unreleased. If the predecessor is released but the current
 source is not, dependence on the current source occurs only through a
 constant gate.
 \item The exit uses a public vector inner product. Its source-specific
 dependence is confined to the same flat gates.
 \item The state coupling vanishes with its full first derivative if the
 current source is unreleased. If that source is released, the predecessor
 is available under the exhaustive release schedule.
\end{itemize}
Thus the entire value and joint gradient of the primal part are computable
exactly from disclosed sources. Multiplying by the public amplitude
$\theta>0$ does not change this conclusion.

Each dual module requires one forward and one transpose action of the
selected branch map per edge; its values follow from the same actions by
inner products. No bound on the dual vectors is needed. The component
replies therefore agree with the actual fixed instance on the common
event. Equivalently, a first mismatching reply is impossible on that event.

For the final, possibly unqueried output, use the public decoder
\[
 \widehat v=\frac{Z(\widehat S_M)}{2R_s},\qquad\|\widehat v\|\le1.
\]
If the output gradient is at most $\theta\tau_*$,
Theorem~\ref{thm:gradient-certificate} implies
\[
 \langle v_M,\widehat v\rangle\ge\frac1{2R_s}>a_*.
\]
Theorem~\ref{thm:source-recovery}, with the dimension enlargement in
\eqref{eq:optimization-dimension}, excludes this on an event of probability
at least $15/16$. Integrating over the independent private seed preserves
the same product-Haar instance law and proves
\eqref{eq:distributional-hardness}.
\end{proof}

The dual maximization in \eqref{eq:value-identity} is exact and
unconstrained. The preceding certificate is consequently a statement about
$\nabla\Phi_0$ itself, with no change of dual domain or of stationarity
criterion.

\section{Parameter calibration in every regime}
\label{app:calibration}

Fix
\begin{equation}
 L_*=10^8,\qquad C_\Delta=31000,\qquad\tau_*=\frac1{20},
 \qquad C_{\mathrm{sc}}=\frac{16L_*C_\Delta}{\tau_*^2}.
 \label{eq:scaling-constant}
\end{equation}
Suppose the normalized family has amplitude $\theta$ and average dual
curvature at least $L_*/\kappa$. Set
\begin{equation}
 s=\frac{4L_*\varepsilon}{L\theta\tau_*},\qquad
 \alpha=\frac{Ls^2}{L_*},\qquad
 f_i^{\mathrm{phys}}(x,y)=\alpha f_i^0(x/s,y/s).
 \label{eq:physical-scaling}
\end{equation}
Every component is then $L$-smooth, average dual curvature is at least
$L/\kappa=\mu$, and $y^*(0)=0$. Moreover,
\begin{align}
 \Phi^{\mathrm{phys}}(0)-\inf\Phi^{\mathrm{phys}}
 &\le\frac{\varepsilon^2}{L}\,
              \frac{C_{\mathrm{sc}}M}{\theta},\label{eq:physical-gap}\\
 \|\nabla\Phi^{\mathrm{phys}}(sx)\|\le4\varepsilon
 &\quad\Longrightarrow\quad
 \|\nabla\Phi_0(x)\|\le\theta\tau_*.
 \label{eq:physical-gradient}
\end{align}
One physical component query is simulated by one normalized component query
and public scaling, including its function value. When
$A\theta\ge2C_{\mathrm{sc}}$, choose
\begin{equation}
 M=\left\lfloor\frac{A\theta}{C_{\mathrm{sc}}}\right\rfloor
 \ge\frac{A\theta}{2C_{\mathrm{sc}}}.
 \label{eq:stage-count}
\end{equation}
This respects the gap budget. The remaining small-stage case is covered
below by the separate component-search lower bound.

\subsection{Large condition numbers}

If $\kappa\ge4L_*\sqrt n$, use the long relay with
\begin{equation}
 N=\left\lfloor\sqrt{\frac{\kappa}{L_*\sqrt n}}\right\rfloor\ge2,
 \qquad m=N-1,\qquad\theta=1.
 \label{eq:long-calibration}
\end{equation}
Its normalized curvature $\beta/N^2$ is at least $L_*/\kappa$.
If exact attainment of the curvature parameter is desired, append to every
component the same independent dual term $-(L_*/\kappa)u^2/2$. Its
maximizer is zero, its value is zero, and it reveals no hidden information.
It preserves the smoothness bound.

For $A\ge2C_{\mathrm{sc}}$, \eqref{eq:stage-count} and the elementary
rounding bound
\[
 N-1\ge\frac14\sqrt{\frac\kappa{L_*\sqrt n}}
\]
give the information cost
\begin{equation}
 \Omega(nM(N-1))=\Omega(An^{3/4}\sqrt\kappa).
 \label{eq:large-lower}
\end{equation}
All suppressed factors depend only on the displayed numerical constants.

\subsection{Small and intermediate condition numbers}

For $L_*\le\kappa\le4L_*\sqrt n$, use the one-edge relay with
\begin{equation}
 \mu_0=\frac{L_*}{\kappa}\in(0,1],\qquad
 \theta=\frac\beta{\mu_0+\beta}
       =\frac\kappa{L_*\sqrt n+\kappa},\qquad m=1.
 \label{eq:short-calibration}
\end{equation}
The physical curvature is exactly $\mu$. When
$A\theta\ge2C_{\mathrm{sc}}$, Theorem~\ref{thm:distributional-hardness}
gives
\begin{equation}
 \Omega(nM)=\Omega(An\theta).
 \label{eq:short-lower}
\end{equation}
To compare this expression across the full overlap, put $r=\kappa/\sqrt n$.
Then
\[
 \frac{\cR(n,\kappa)}{n\theta}
 =\begin{cases}
   L_*+r,&r\le1,\\[2pt]
   (L_*+r)/\sqrt r,&1\le r\le4L_*.
  \end{cases}
\]
Both expressions are at most $5L_*$. Hence
\begin{equation}
 \cR(n,\kappa)\le5L_*n\theta
 \qquad(L_*\le\kappa\le4L_*\sqrt n).
 \label{eq:short-comparison}
\end{equation}
This gives the desired $\Omega(A\cR)$ term whenever stages fit.
If $A\theta<2C_{\mathrm{sc}}$, instead
\begin{equation}
 A\cR(n,\kappa)<10L_*C_{\mathrm{sc}}n,
 \label{eq:no-stage}
\end{equation}
so the explicit $\Omega(n)$ construction in
Section~\ref{sec:search-lower} covers this entire case. No extra
requirement such as $A\gtrsim\sqrt n/\kappa$ enters the main theorem.

Finally, if $1\le\kappa<L_*$, then
$\cR(n,\kappa)\le L_*\sqrt n$. The unrestricted minimization baseline in
Section~\ref{sec:minimization-baseline}, embedded with a decoupled dual
quadratic, gives the required bound throughout this bounded range.

\section{Stopping times, baseline instances, and upper-bound interfaces}
\label{app:protocols}

\subsection{From distributional hardness to expected query cost}
\label{sec:expected-lower}

Suppose an algorithm has expected query count at most $Q$ on every physical
instance and satisfies
$\E\|\nabla\Phi(\widehat x)\|^2\le\varepsilon^2$.
Under the fixed hard distribution, Markov's inequality gives
\[
 \PP\bigl(\|\nabla\Phi(\widehat x)\|\le4\varepsilon\bigr)
 \ge\frac{15}{16}.
\]
Let $b=c_{\mathrm{inf}}nMm$ and $B=\lfloor b\rfloor$. If $B\ge1$,
truncate the algorithm after $B$ calls, using a fixed fallback output on
overflow. Its success probability at threshold $4\varepsilon$ is at least
$15/16-Q/B$. By Theorem~\ref{thm:distributional-hardness} and
\eqref{eq:physical-gradient}, it is at most $1/16$. Therefore
\[
 Q\ge\frac78B\ge\frac7{16}b.
\]
If $b<1$, the same hard distribution defeats every zero-query output.
Compare the algorithm with its zero-query truncation. Their outputs can
differ only if at least one query is made, an event of probability at most
the expected number of calls, hence at most $Q$. It follows that
$Q\ge7/8\ge7b/8$. Thus all cases yield $\Omega(nMm)$ without a rounding
gap. A worst-case expected guarantee holds on each instance, so it also
holds after averaging over the fixed hard distribution.

The fixed-budget constant-confidence lower bound follows directly from
Theorem~\ref{thm:distributional-hardness}: success at threshold
$\varepsilon$ implies success at threshold $4\varepsilon$. An integer
budget exceeding $\lfloor b\rfloor$ is strictly larger than $b$.

\subsection{An explicit additive component-search lower bound}
\label{sec:search-lower}

Let $J$ be uniform on $[n]$, and let $s_0$ be an independent uniform sign.
With scalar primal and dual variables, define
\begin{equation}
 f_i(x,y)=\frac L2x^2-nb_0s_0\mathbf1_{\{i=J\}}x-\frac\mu2y^2,
 \qquad b_0=8\varepsilon.
 \label{eq:search-instance}
\end{equation}
Every component is $L$-smooth, $y^*(0)=0$, and
\[
 \Phi(x)=\frac L2x^2-b_0s_0x,\qquad
 \Phi(0)-\inf\Phi=\frac{32\varepsilon^2}{L}\le\Delta
\]
whenever $A\ge32$. Before a query to $J$, values and gradients are
independent of $s_0$. A $q$-query algorithm hits $J$ with probability at
most $q/n$. On the no-hit event, an output cannot satisfy
$|Lx-b_0s_0|\le4\varepsilon$ for both signs. Its success probability is
at most
\begin{equation}
 \frac12+\frac{q}{2n}.
 \label{eq:search-success}
\end{equation}
For expected-cost algorithms, take $q=\lfloor n/4\rfloor$ when $n\ge4$.
The truncation comparison gives
$15/16-Q/q\le5/8$, hence $Q\ge5q/16=\Omega(n)$.
For $n<4$, zero-query success is at most $1/2$, which forces
$Q\ge7/16=\Omega(n)$ by the same comparison with a zero-query
truncation. For deterministic budgets, success probability $2/3$ in
\eqref{eq:search-success} requires at least $n/3$ queries. This proves the
additive term under both performance criteria.

\subsection{The imported minimization baseline}
\label{sec:minimization-baseline}

We use the individually smooth nonconvex finite-sum lower bound of
\citet[Theorem 3.1 and Appendices F.3 and G.1]{pengdense2026}.
For a universal accuracy condition and all $n\ge n_0$, it gives
$\Omega(n+A\sqrt n)$ against unrestricted randomized value-and-gradient
queries, including worst-case expected query cost with expected gradient
norm and fixed-budget success probability $2/3$.
The specified failure constant is $11/16$, so the latter confidence is
covered. Expected squared-gradient accuracy implies expected-norm accuracy
by Jensen's inequality.

For a hard minimization finite sum $n^{-1}\sum_i g_i(x)$, set
\begin{equation}
 f_i(x,y)=g_i(x)-\frac\mu2\|y\|^2.
 \label{eq:baseline-embedding}
\end{equation}
The embedding is $L$-smooth and exactly initialized, leaves the primal
value unchanged, and uses one minimization query per minimax query.

For $n<n_0$, fix $N_0=n_0$ and an $N_0$-component hard average $G$.
Make every target component $G(x)-\mu\|y\|^2/2$; each target query costs
$N_0$ source queries. The inherited bound
$\Omega(1+A/\sqrt{N_0})$ is $\Omega(n+A\sqrt n)$ after changing universal
constants, since $1\le n<N_0$ and $N_0$ is fixed.

\subsection{Combining the lower bounds}
\label{sec:combine-lower}

Choose $A_*$ larger than $2C_{\mathrm{sc}}$, $32$, and the reciprocals of
the universal accuracy thresholds in the imported baseline and upper
algorithms. There are three cases.
\begin{enumerate}[label=(\roman*)]
 \item If $\kappa\ge4L_*\sqrt n$, use \eqref{eq:large-lower}.
 \item If $L_*\le\kappa\le4L_*\sqrt n$, use
 \eqref{eq:short-lower}--\eqref{eq:short-comparison} when stages fit, and
 \eqref{eq:no-stage} with the search lower bound otherwise.
 \item If $1\le\kappa<L_*$, use the minimization baseline and
 $\cR(n,\kappa)\le L_*\sqrt n$.
\end{enumerate}
The separate $n$ term is always available. Taking the maximum of these
hard families and using $\max\{u,v\}\ge(u+v)/2$ gives
\[
 \T^{\mathrm E}_{\mathrm{IS}}\ge c[n+A\cR(n,\kappa)]
\]
for universal $c>0$, under both performance protocols. All instances satisfy
IS and hence AS. The factor $nMm$ comes from the serial information theorem.

\subsection{Imported AS upper bounds and fixed-confidence truncation}
\label{sec:upper-interfaces}

Under AS, average dual strong concavity, and \eqref{eq:initialization},
RPDA has fixed cost $\widetilde O(n+A\sqrt n\,\kappa)$ for $n\ge2$ and
$\kappa\le\sqrt n$ \citep[Theorem 3 and Appendix H]{zhou2026}.
For $\kappa\ge\sqrt n$, Catalyst--SVRG has expected cost
$\widetilde O(A(n+n^{3/4}\sqrt\kappa))$
\citep[Theorem 4.1, Corollary 4.1, and Section 4.2]{zhang2021}.
Both guarantee $\E\|\nabla\Phi(\widehat x)\|^2\le\varepsilon^2$ using
only gradients; the initial dual-error term is zero.
Since $n^{3/4}\sqrt\kappa\ge n$ in the second regime and $A\ge1$, the
combined cost is $\widetilde O(n+A\cR)$.
For $n=1$, simulate the $n=2$ algorithm on two identical copies,
at one actual call per simulated component call.

For a deterministic budget, run at accuracy $\varepsilon/4$ and truncate
after $\lceil12Q_0\rceil$ calls, where $Q_0$ is the public expected-cost
bound. Markov's inequality bounds overflow by $1/12$ and the untruncated
gradient error probability by $1/16$, giving success at least
$1-1/12-1/16=41/48>2/3$. The accuracy change multiplies $A$ by $16$
and leaves $\kappa$ unchanged. This proves Theorem~\ref{thm:main}; no
expected-square claim is made for the fallback output.


\end{document}
